% !TEX program = lualatex
% =============================================================================
%  Guia Rápido de Uso — Hanaro — modelo AX Academy / IFAM
%
%  COMPILAR NA RAIZ: uv run build hanaro (LuaLaTeX, três passadas)
%
%  Todo o estilo mora em axacademy-guia.sty. Aqui ficam só os dados do guia e
%  a ordem dos capítulos, que estão em caps/. As telas ficam em img/.
%
%  O conteúdo é o do Hanaro, o sistema de controle de perdas (scrap) da Linha
%  Final, e serve de gabarito para os outros projetos: telas reais do MVP,
%  Os capítulos 1 e 2 conservam os oito passos originais. Os demais ampliam a consulta.
%
%  Sem número de equipe em nenhum lugar do guia: a autoria é a equipe do Polo
%  INOVA. Subtítulo e tagline vazios tiram o texto da capa, da folha de rosto
%  e da contracapa.
% =============================================================================
\documentclass[11pt,a4paper,oneside,openany]{book}
\usepackage{axacademy-guia}

% ---------- DADOS DESTE GUIA ----------
\guiatitulo{Guia do\\Usuário}
\guiasubtitulo{}
\guiasistema{Hanaro}
\guiatagline{}
\guiaequipe{equipe do Polo INOVA}
\guiaano{2026}
\guiaversao{1.6}
\guiacidade{Manaus, Amazonas --- Brasil}
\guiacliente{LG Electronics do Brasil}
\guiareferencia{}

% ---------- CONTATO (folha de rosto) ----------
% Campo vazio some da folha de rosto, com o ícone. O suporte do Hanaro é feito
% pelo próprio sistema (Capítulo 10), por isso não há e-mail.
\guiaemail{}
\guiasite{www.ifam.edu.br}
\guiaendereco{INOVA --- Polo de Inovação do IFAM\\
    Av. Gov. Danilo de Matos Areosa, 1975 --- Distrito Industrial I\\
    Manaus/AM --- CEP 69075-351}

% ---------- AVISO DA CAPA ----------
% Quebras explícitas: a caixa comporta três linhas curtas.
\guiaaviso{LEIA ANTES DO\\PRIMEIRO ACESSO}

% ---------- MARCA DO SISTEMA ----------
% O Hanaro tem o logotipo inteiro em arquivo (art/marca/hanaro-assinatura.pdf),
% então o ícone fica vazio e o nome não é composto em texto.
\guiamarcaicone{}
\guiamarcatexto{hanaro-assinatura}

% ---------- IMAGEM DA CAPA ----------
% Vazio: capa sem imagem. Para pôr uma, nomeie um arquivo de img/.
\guiacapaimagem{}

% ---------- TEXTO DA CONTRACAPA ----------
% Sai logo abaixo das marcas da AX Academy e do IFAM.
\guiasobreaequipe{Orientações de uso e consulta operacional — novembro de 2026}

\begin{document}

% ---------- ABERTURA ----------
\guiacapa
\guiafolhaderosto
\guiaavisodeuso

% ---------- HISTÓRICO DE REVISÕES ----------
\begin{guiahistorico}
    \guiarevisao{1.0}{Outubro de 2026}{AX Academy}{Emissão inicial.}
    \guiarevisao{1.2}{Novembro de 2026}{Equipe do Polo INOVA}{Preservação dos capítulos 1 e 2 do guia rápido e inclusão da consulta detalhada a partir do capítulo 3.}
    \guiarevisao{1.3}{Novembro de 2026}{Equipe do Polo INOVA}{Consolidação de passos repetidos e referências ao guia rápido, com os capítulos 1 e 2 preservados.}
    \guiarevisao{1.4}{Novembro de 2026}{Equipe do Polo INOVA}{Revisão dos marcadores, recortes dos formulários e correspondência entre imagem e instruções.}
    \guiarevisao{1.5}{Novembro de 2026}{Equipe do Polo INOVA}{Revisão editorial, padronização das instruções e esclarecimento das regras e dos indicadores.}
    \guiarevisao{1.6}{Novembro de 2026}{Equipe do Polo INOVA}{Consolidação de Configurações e IA, revisão de permissões e orientações conforme a versão das capturas.}
\end{guiahistorico}

\begin{remark}[info]
Os capítulos 1 e 2 reproduzem integralmente o guia rápido do perfil Gestor. As indicações de público e escopo nesses capítulos se referem ao uso básico. A partir do capítulo 3, este mesmo volume apresenta a consulta detalhada para os perfis Analista, Gestor e Administrador.
\end{remark}

\clearpage

\guiasobreax

\guiapagina{Sobre o sistema}{%
    O Hanaro é o sistema web de controle das perdas produtivas (\emph{scrap}) da
    Linha Final. Ele recebe o relatório \emph{Other Account Transaction Text}
    exportado do GERP, consolida a base de scrap e publica o IF Cost no Dashboard e
    no Modo TV, exibido nas telas da fábrica.

    Antes dele, o analista de Produção reprocessava à mão a planilha enviada pela
    Qualidade, de 1,5 a 2 horas por dia, para converter os valores, recortar o
    escopo e eleger os ofensores do período.

    O fluxo básico tem oito passos: entrar no sistema, enviar o relatório do GERP,
    conferir a execução, ler o Dashboard, revisar a Base de Scrap, consultar a
    análise de uma ocorrência, aplicar um modelo de análise em lote e exibir o Modo
    TV. A \autoref{sec:fluxo} apresenta esse fluxo em BPMN; o \autoref{cap:passos}
    descreve cada passo com a tela correspondente.
}

% Espaçamento do índice ajustado ao volume completo; conteúdo dos capítulos preservado.
\begingroup
\patchcmd{\guiasumario}{\setlength{\cftbeforechapskip}{8pt}}{\setlength{\cftbeforechapskip}{5pt}}{}{\PackageError{hanaro}{Espacamento de capitulos do indice nao localizado}{Verifique o estilo do guia.}}
\patchcmd{\guiasumario}{\setlength{\cftbeforesecskip}{2pt}}{\setlength{\cftbeforesecskip}{1pt}}{}{\PackageError{hanaro}{Espacamento de secoes do indice nao localizado}{Verifique o estilo do guia.}}
\guiasumario
\endgroup

% ---------- PARTE I ----------
% Um capítulo só por parte; dentro dele, só seções.
\part{Visão Geral}
% =============================================================================
%  Capítulo 1 — Antes de Começar (Parte I, Visão Geral)
%  Capítulo único da parte: público, requisitos de acesso, como ler o guia e
%  o fluxo de uso em BPMN, com a legenda e o mapa passo -> tela -> página.
% =============================================================================
\chapter{Antes de Começar}

\section{Público e escopo}

Este guia se destina ao usuário com perfil \textbf{Gestor}, que mantém a base de
scrap atualizada e acompanha o IF Cost da Linha Final. A leitura completa leva cerca
de quinze minutos.

\begin{itemize}
    \item \textbf{Coberto:} acesso, envio do relatório do GERP, conferência da
          execução, Dashboard, Base de Scrap, análise de uma ocorrência, aplicação de
          modelo em lote e Modo TV.
    \item \textbf{Não coberto:} relatórios de análise, alertas e regras de
          notificação, configurações (classificações de material, tipos de scrap,
          programa do ano) e administração de usuários.
\end{itemize}

\section{Requisitos de acesso}

\begingroup
\guiazebra
\small
\noindent\begin{longtable}{@{}L{4.0cm}L{11.5cm}@{}}
    \toprule
    \guiacabtabela
    \thead{Item} & \thead{Requisito} \\
    \midrule
    \endhead
    \bottomrule
    \endlastfoot
    Navegador           & Navegador atualizado. As telas deste guia foram capturadas no
                          Google Chrome \\
    Endereço do sistema & Informado pela equipe na liberação do acesso \\
    Credenciais         & Usuário e senha do Hanaro, cadastrados pelo administrador \\
    Perfil              & Gestor \\
    Arquivo de entrada  & Relatório \emph{Other Account Transaction Text} exportado do
                          GERP, em texto separado por tabulação (TSV), com codificação
                          UTF-8 ou CP1252 \\
    Modo TV             & TV ou monitor da linha com navegador, para a exibição em tela
                          cheia \\
\end{longtable}
\endgroup

\begin{remark}
O usuário só entra no Hanaro depois de cadastrado e ativado pelo administrador do
sistema. Sem esse cadastro, a tela de acesso recusa as credenciais.
\end{remark}

\section{Como ler este guia}

O \autoref{cap:passos} descreve o uso em \textbf{passos}. Cada passo ocupa uma linha
da página, dividida em duas colunas:

\begin{itemize}
    \item \textbf{à esquerda, a tela} do sistema, exatamente como o usuário a vê;
    \item \textbf{à direita, o texto} do passo: o que fazer, em que ordem e o que o
          sistema responde.
\end{itemize}

Os números em vermelho sobre a tela, como \marca{1} e \marca{2}, indicam os elementos
citados no texto ao lado, na mesma ordem. O número do passo, no selo
\colorbox{guiavermelho}{\color{white}\guiahead\bfseries\fontsize{7.5}{7.5}\selectfont PASSO 1},
é o mesmo que aparece no diagrama BPMN da \autoref{sec:fluxo}.

Botões e opções da interface aparecem em destaque, como \botao{Entrar}. Um caminho de
menu é escrito como \caminho{Menu\seta Execuções}.

As caixas ao longo do texto têm três tipos:

\begin{remark}
\textbf{Atenção.} Condição que impede o passo de ser concluído ou que exige
correção.
\end{remark}

\begin{remark}[dica]
\textbf{Dica.} Atalho ou forma mais rápida de concluir o passo.
\end{remark}

\begin{remark}[info]
\textbf{Informação.} Comportamento do sistema que convém conhecer, sem ação
necessária.
\end{remark}

% ---------------------------------------------------------------------------
%  Fluxo de uso em BPMN
%
%  O diagrama usa os estilos `bpmn ...' do axacademy-guia.sty. As raias são
%  COLUNAS e o fluxo desce: assim o diagrama cabe em pé, nos 16 cm da largura
%  do texto. Coordenadas em cm, com y = 0 no topo das raias.
%
%  A seção abre página nova: o diagrama vai em `figure[H]', que não flutua, e
%  precisa da página inteira. Com oito passos, as tarefas têm uma linha só e
%  as linhas do diagrama ficam a 1,15 cm uma da outra; com mais passos,
%  divida o fluxo em dois diagramas.
% ---------------------------------------------------------------------------
\clearpage
\section{Fluxo de uso}\label{sec:fluxo}

O fluxo básico vai do acesso ao Modo TV. A raia \textbf{Gestor} reúne o que o
usuário faz; a raia \textbf{Hanaro}, o que o sistema faz em resposta. As tarefas com
borda vermelha são os passos deste guia: o número no canto é o do passo e leva à
página em que ele é descrito.

\begin{figure}[H]
\centering
\begin{tikzpicture}[x=1cm,y=1cm,
        bpmn tarefa/.append style={text width=4.4cm,minimum height=0.72cm}]
    % A caixa de medida é a das raias, exata: sem ela, a espessura do contorno
    % passa 0,7 pt da largura do texto.
    \useasboundingbox (0,0) rectangle (16,-13.35);
    % ---- raias ----
    \bpmnraia{0}{8}{13.35}{Gestor}
    \bpmnraia{8}{8}{13.35}{Hanaro}

    % ---- gestor ----
    \node[bpmn inicio] (ini) at (4,-1.05) {};
    \node[bpmn rotulo,anchor=west] at ([xshift=4pt]ini.east)
        {Relatório do GERP\\disponível};
    \bpmnpasso{p1}{4,-2.15}{1}{Entrar com usuário e senha}
    \bpmnpasso{p2}{4,-4.45}{2}{Enviar o relatório do GERP}
    \bpmnpasso{p3}{4,-6.75}{3}{Conferir a execução}
    \bpmnpasso{p4}{4,-7.90}{4}{Ler o Dashboard}
    \bpmnpasso{p5}{4,-9.05}{5}{Revisar a Base de Scrap}
    \bpmnpasso{p6}{4,-10.20}{6}{Consultar a análise do caso}
    \bpmnpasso{p7}{4,-11.35}{7}{Aplicar o modelo em lote}
    \bpmnpasso{p8}{4,-12.50}{8}{Exibir o Modo TV}

    % ---- sistema ----
    \node[bpmn tarefa] (s1) at (12,-2.15) {Validar as credenciais};
    \bpmngateway{g1}{12,-3.30}
    \node[bpmn rotulo,anchor=west] at ([xshift=3pt]g1.east) {Credenciais\\válidas?};
    \node[bpmn tarefa] (s2) at (12,-4.45) {Abrir o Dashboard};
    \node[bpmn tarefa] (s3) at (12,-5.60) {Processar o arquivo enviado};
    \node[bpmn tarefa] (s4) at (12,-6.75) {Publicar o snapshot da base};
    \node[bpmn tarefa] (s5) at (11.6,-12.50) {Girar as telas do Modo TV};
    \node[bpmn fim] (fim) at (14.95,-12.50) {};
    \node[bpmn rotulo,anchor=south,align=center] at ([yshift=3pt]fim.north)
        {Indicadores\\na fábrica};

    % ---- fluxos ----
    \draw[bpmn fluxo] (ini) -- (p1);
    \draw[bpmn fluxo] (p1) -- (s1);
    \draw[bpmn fluxo] (s1) -- (g1);
    \draw[bpmn fluxo] (g1) -- node[bpmn rotulo,anchor=west,xshift=2pt] {Sim} (s2);
    \draw[bpmn fluxo] (g1.west) -- node[bpmn rotulo,anchor=south,pos=0.25] {Não}
        (4,-3.30) -- (p1.south);
    \draw[bpmn fluxo] (s2) -- (p2);
    \draw[bpmn fluxo] (p2.south) |- (s3.west);
    \draw[bpmn fluxo] (s3) -- (s4);
    \draw[bpmn fluxo] (s4) -- (p3);
    \draw[bpmn fluxo] (p3) -- (p4);
    \draw[bpmn fluxo] (p4) -- (p5);
    \draw[bpmn fluxo] (p5) -- (p6);
    \draw[bpmn fluxo] (p6) -- (p7);
    \draw[bpmn fluxo] (p7) -- (p8);
    \draw[bpmn fluxo] (p8) -- (s5);
    \draw[bpmn fluxo] (s5) -- (fim);
\end{tikzpicture}
\caption{Fluxo básico de uso do Hanaro pelo Gestor, em BPMN. Os números indicam os
passos do \autoref{cap:passos}.}
\label{fig:bpmn}
\end{figure}

\section{Leitura do diagrama}

\begingroup
\small
\renewcommand{\arraystretch}{1.6}
\noindent\begin{tabularx}{\textwidth}{@{}M{2.4cm}X@{}}
    \toprule
    \guiacabtabela
    \thead{Símbolo} & \thead{Significado} \\
    \midrule
    \tikz[baseline=-0.6ex]\node[bpmn inicio,minimum size=0.42cm]{};
        & Início do fluxo: o evento que leva o usuário ao sistema. \\
    \tikz[baseline=-0.6ex]\node[bpmn tarefa,text width=1.2cm,minimum height=0.5cm,
            font=\fontsize{7}{8}\selectfont]{Tarefa};
        & Tarefa executada pelo usuário ou pelo sistema, conforme a raia em que está. \\
    \tikz[baseline=-0.6ex]{\node[bpmn tarefa,draw=guiavermelho,line width=1.1pt,
            text width=1.2cm,minimum height=0.5cm,font=\fontsize{7}{8}\selectfont] (t) {Passo};
            \node[guia marcador,minimum size=11pt,font=\guiahead\bfseries\fontsize{7}{7}\selectfont]
            at (t.north west) {3};}
        & Tarefa descrita neste guia. O número é o do passo; no PDF, ele leva à página
          correspondente. \\
    \tikz[baseline=-0.6ex]{\bpmngateway{g}{0,0}}
        & Decisão. O fluxo segue por um único caminho, conforme a condição escrita ao
          lado. \\
    \tikz[baseline=-0.6ex]\draw[bpmn fluxo] (0,0) -- (1.2,0);
        & Sequência: a próxima tarefa só começa quando a anterior termina. \\
    \tikz[baseline=-0.6ex]\node[bpmn fim,minimum size=0.42cm]{};
        & Fim do fluxo. \\
    \bottomrule
\end{tabularx}
\endgroup

\section{Mapa dos passos}

\begingroup
\guiazebra
\small
\noindent\begin{longtable}{@{}P{1.3cm}L{3.2cm}L{8.4cm}P{1.4cm}@{}}
    \toprule
    \guiacabtabela
    \thead{Passo} & \thead{Tela} & \thead{Ação principal} & \thead{Página} \\
    \midrule
    \endhead
    \bottomrule
    \endlastfoot
    \marca{1} & Acesso           & Informar usuário e senha e clicar em \botao{Entrar}       & \pagina{passo:1} \\
    \marca{2} & Execuções        & Selecionar o arquivo do GERP e clicar em
                                   \botao{Processar relatório}                               & \pagina{passo:2} \\
    \marca{3} & Execuções        & Conferir o status e o snapshot da execução                & \pagina{passo:3} \\
    \marca{4} & Dashboard        & Ler o IF Cost acumulado, o target e os rankings           & \pagina{passo:4} \\
    \marca{5} & Base de Scrap    & Filtrar e revisar as ocorrências de scrap                 & \pagina{passo:5} \\
    \marca{6} & Análise de scrap & Ler o custo, a contenção e as evidências da ocorrência    & \pagina{passo:6} \\
    \marca{7} & Base de Scrap    & Marcar as ocorrências e aplicar o modelo de análise       & \pagina{passo:7} \\
    \marca{8} & Modo TV          & Abrir o Modo TV e controlar a rotação das telas           & \pagina{passo:8} \\
\end{longtable}
\endgroup

\fimdecapitulo


% ---------- PARTE II ----------
\part{Passo a Passo}
% =============================================================================
%  Capítulo 2 — Uso Básico do Hanaro (Parte II, Passo a Passo)
%  Capítulo único da parte, só em seções: os oito passos, as colunas do
%  histórico de execuções, as telas do Modo TV e o encerramento da sessão.
%
%  Cada \passo é uma linha de duas colunas: a tela à esquerda e o texto à
%  direita, e conta como uma seção do capítulo (2.1, 2.2...). O argumento
%  opcional lista as posições dos marcadores sobre a tela, como x/y em fração
%  da imagem, com origem no canto SUPERIOR esquerdo. A numeração segue a
%  ordem da lista e casa com a lista `elementos'.
%
%      \passo[0.335/0.525, 0.665/0.649]{01-login}{Título}{Texto ...}
% =============================================================================
\chapter{Uso Básico do Hanaro}\label{cap:passos}

% ---------- acesso e atualização da base ----------
\passo[0.335/0.525, 0.665/0.649, 0.335/0.772]{01-login}{Entrar no Hanaro}{%
    Abra o endereço do Hanaro no navegador. A tela de acesso é a primeira a aparecer.
    \begin{elementos}
        \item \campo{Usuário}: o nome de usuário cadastrado pelo administrador. Não é
              o e-mail.
        \item \campo{Senha}: a senha do usuário. O ícone de olho, no fim do campo,
              mostra o que foi digitado.
        \item \botao{Entrar}: valida as credenciais e abre o Dashboard.
    \end{elementos}
    Credenciais inválidas mantêm o usuário nesta tela.
}

\passo[0.115/0.181, 0.845/0.146, 0.31/0.521, 0.69/0.669]{02-enviar-relatorio}{Enviar o relatório do GERP}{%
    A base de scrap é atualizada pelo relatório \emph{Other Account Transaction
    Text}. Exporte-o do GERP antes deste passo.
    \begin{elementos}
        \item \caminho{Execuções}, no menu lateral: abre a tela
              \textbf{Execuções e atualização}.
        \item \botao{Enviar relatório}: abre a janela \textbf{Atualizar base por
              relatório}.
        \item \textbf{Área do arquivo}: um clique abre a escolha do arquivo exportado
              do GERP. O nome do arquivo escolhido aparece no lugar de
              \emph{Selecionar relatório GERP}.
        \item \botao{Processar relatório}: envia o arquivo. O Hanaro registra uma
              nova execução, conferida no passo seguinte.
    \end{elementos}
    \botao{Cancelar} fecha a janela sem enviar.
}

\begin{remark}
O arquivo precisa ser o \emph{Other Account Transaction Text}, em texto separado por
tabulação (TSV), com codificação UTF-8 ou CP1252, como indica a própria janela.
Planilha salva em outro formato não serve.
\end{remark}

\passo[0.148/0.378, 0.98/0.378, 0.871/0.242]{03-execucoes}{Conferir a execução}{%
    A execução nova entra no topo do \textbf{Histórico de rotinas}.
    \begin{elementos}
        \item \textbf{Linha da execução}: o identificador (\texttt{EXE-...}), a
              origem e o gatilho. Um clique na linha abre os \textbf{Detalhes da
              execução}, com o arquivo de origem e as contagens.
        \item \textbf{Status}: \emph{Na fila} enquanto a execução aguarda o
              processamento e \emph{Concluído} ao final. Ao lado, o
              \textbf{Snapshot} passa de \emph{Não publicado} a \emph{Publicado}.
        \item \botao{Filtrar}: restringe o histórico.
    \end{elementos}
    As colunas estão descritas na \autoref{tab:historico}.
}

\begin{remark}[info]
O Dashboard passa a refletir o arquivo novo quando o snapshot da execução fica
\emph{Publicado}. Até lá, ele mostra a base da execução anterior.
\end{remark}

\section{Colunas do histórico}

% Tabela curta: vai inteira, em `table', para não partir entre páginas.
\begin{table}[H]
\caption{Colunas do histórico de rotinas.}\label{tab:historico}
\guiazebra
\small
\noindent\begin{tabularx}{\textwidth}{@{}L{4.4cm}X@{}}
    \toprule
    \guiacabtabela
    \thead{Coluna} & \thead{Conteúdo} \\
    \midrule
    \textbf{Execution ID / Processo} & Identificador da execução e nome do relatório
                                       processado \\
    \textbf{Origem}                  & Sistema de onde veio o arquivo: GERP \\
    \textbf{Gatilho}                 & Como a execução começou. \emph{Manual}: pelo
                                       envio do relatório no passo~2 \\
    \textbf{Início}, \textbf{Fim} e \textbf{Duração}
                                     & Data e hora do processamento e o tempo que ele
                                       levou \\
    \textbf{Recebidos}               & Registros lidos no arquivo \\
    \textbf{Válidos}                 & Registros aceitos na base de scrap \\
    \textbf{Rejeitados}              & Registros recusados na validação \\
    \textbf{Snapshot}                & \emph{Publicado} quando a base processada já
                                       alimenta os indicadores; \emph{Não publicado}
                                       antes disso \\
    \textbf{Status}                  & \emph{Na fila} enquanto aguarda;
                                       \emph{Concluído} ao final \\
    \bottomrule
\end{tabularx}
\end{table}

\begin{remark}
Com \textbf{Rejeitados} acima de zero, parte do arquivo ficou fora da base. Abra os
detalhes da execução e acione o suporte em \caminho{Ajuda e suporte}, no pé do menu
lateral, com o identificador da execução.
\end{remark}

% ---------- indicadores e ocorrências ----------
\passo[0.30/0.114, 0.36/0.214, 0.66/0.114, 0.45/0.322, 0.924/0.503, 0.27/0.725]{04-dashboard}{Ler o Dashboard}{%
    O Dashboard é a tela que abre depois do acesso. No alto, à direita,
    \emph{Atualizado às} indica a hora da última atualização.
    \begin{elementos}
        \item \textbf{Impacto absoluto}: chave que alterna os indicadores entre o IF
              Cost em dólar e o Scrap Rate, em percentual.
        \item \textbf{Indicadores}: IF Cost acumulado, target, atingimento do target,
              diferença para a meta e variação sobre o mesmo acumulado do ano
              anterior.
        \item \textbf{Controles}: o botão com o ícone de tela abre o Modo TV
              (passo~8); \botao{YTD anterior} escolhe a base de comparação;
              \botao{Filtros} recorta os dados; \textbf{Valor monetário} alterna
              entre dólar e quantidade, em unidades.
        \item \textbf{Target × realizado mensal}: realizado do ano, referência do ano
              anterior e target, mês a mês.
        \item \textbf{Top 10 produtos / áreas}: o IF Cost por produto ou área.
        \item \textbf{Top 10 linhas} e \textbf{Top 10 modelos}: os rankings de
              ofensores. Mais abaixo, componentes e ofensores.
    \end{elementos}
}

\begin{remark}[dica]
Cada quadro tem o seu botão \botao{Filtros} e um botão de tabela, que abre os mesmos
números em linhas e colunas. Passar o cursor sobre um ponto do gráfico mostra os
valores do mês.
\end{remark}

\passo[0.12/0.178, 0.555/0.194, 0.135/0.337, 0.70/0.1]{05-base-scrap}{Revisar a Base de Scrap}{%
    A lista traz as ocorrências consolidadas pelas execuções. Ao lado do título, o
    total de ocorrências, o valor em dólar e a quantidade.
    \begin{elementos}
        \item \caminho{Base de Scrap}, no menu lateral: abre a tela.
        \item \textbf{Ordem e filtros}: \botao{Data da transação} escolhe o campo da
              ordem e \botao{Decrescente} a inverte; a lista \botao{Todos} filtra
              pelo status da análise; \botao{Filtrar} abre os demais filtros.
        \item \textbf{Ocorrência}: organização e produto, data, item, descrição,
              ordem, quantidade, valor, tipo de scrap, status da análise
              (\emph{Não analisado} ou \emph{Revisado}), responsável e relatório. Um
              clique abre a análise (passo~6).
        \item \botao{Aplicar Modelo} e \botao{Selecionar itens}: levam uma análise já
              feita a várias ocorrências (passo~7).
    \end{elementos}
}

\passo[0.39/0.138, 0.195/0.349, 0.195/0.507, 0.195/0.652, 0.577/0.87]{06-ocorrencia}{Consultar a análise de uma ocorrência}{%
    A janela \textbf{Análise de scrap} reúne o que se sabe da ocorrência. Role a
    janela para ver todas as seções.
    \begin{elementos}
        \item \textbf{Status da análise}: \emph{Revisado}, como nesta ocorrência, ou
              \emph{Não analisado}.
        \item \textbf{Custo}: tipo de custo, modelo e quantidade. Acima, \textbf{Onde
              e quando} traz a data e a hora.
        \item \textbf{Contenção imediata}: a ação tomada na linha para conter o
              problema.
        \item \textbf{Evidências e fotos}: as imagens anexadas à análise.
        \item \botao{Criar relatório desta análise}: abre um relatório de análise a
              partir desta ocorrência.
    \end{elementos}
    O coração, no alto, guarda a análise como modelo para o passo seguinte.
}

\passo[0.565/0.09, 0.145/0.494, 0.313/0.632, 0.685/0.768]{07-modelo-lote}{Aplicar um modelo em lote}{%
    Quando o mesmo problema se repete em várias ocorrências, a análise de uma serve
    de modelo para as outras.
    \begin{elementos}
        \item \textbf{Modelo em uso}: em \botao{Aplicar Modelo}, escolha o modelo em
              \textbf{Modelos salvos} e clique em \botao{Usar este modelo}. A faixa
              no alto confirma a escolha.
        \item \textbf{Ocorrências marcadas}: com \botao{Selecionar itens}, marque as
              ocorrências e clique em \botao{Aplicar modelo aos itens selecionados},
              acima da lista. Abre esta janela.
        \item \textbf{Reaproveitar imagens do relatório de referência}: copia as
              fotos do modelo para todas as ocorrências marcadas.
        \item \botao{Aplicar modelo aos itens selecionados}: grava a análise do
              modelo nas ocorrências marcadas.
    \end{elementos}
}

\begin{remark}
Ocorrências inativas ou que já têm relatório são ignoradas na aplicação, como avisa a
própria janela. Copiar as fotos replica as mesmas evidências em todas as ocorrências
marcadas: use só quando as fotos valem para todas.
\end{remark}

% ---------- Modo TV ----------
\passo[0.395/0.069, 0.712/0.083, 0.2/0.865, 0.5/0.865, 0.785/0.865]{08-modo-tv}{Exibir o Modo TV}{%
    O Modo TV mostra os indicadores em tela cheia, para as TVs da fábrica. Ele abre
    pelo botão com o ícone de tela, nos controles do Dashboard (passo~4).
    \begin{elementos}
        \item \textbf{Tela em exibição}: o título da tela. Na tela Fábrica,
              \botao{Mês} e \botao{Ano} escolhem o período.
        \item \textbf{Dados até}: a data do último dado da base e, ao lado, a data e
              a hora atuais.
        \item \textbf{Telas}: \botao{Executivo}, \botao{Fábrica}, \botao{Ofensores} e
              \botao{Prioridades}, descritas na \autoref{tab:telastv}. Um clique leva
              direto à tela.
        \item \textbf{Rotação}: as setas passam de tela, \botao{Pausar} para a
              rotação automática e \textbf{Intervalo} define o tempo de cada tela.
        \item \botao{Tema Claro}, \botao{Tela cheia} e \botao{Sair}: trocam o fundo
              escuro pelo claro, ocupam a tela inteira e voltam ao Dashboard.
    \end{elementos}
}

\begin{remark}[dica]
Na TV da linha, deixe a \textbf{Rotação automática} ligada e use \botao{Tela cheia}:
as quatro telas se alternam sem intervenção. O contador no fim da barra, como
\emph{2/4}, indica a tela atual.
\end{remark}

\section{Telas do Modo TV}

% Tabela curta: vai inteira, em `table', para não partir entre páginas.
\begin{table}[H]
\caption{As quatro telas do Modo TV.}\label{tab:telastv}
\guiazebra
\small
\noindent\begin{tabularx}{\textwidth}{@{}L{2.4cm}L{3.9cm}X@{}}
    \toprule
    \guiacabtabela
    \thead{Botão} & \thead{Título da tela} & \thead{O que mostra} \\
    \midrule
    \textbf{Executivo}   & Visão Executiva
                         & IF Cost acumulado, mesmo período do ano anterior, variação e
                           target acumulado; evolução do IF Cost no ano, no ano anterior
                           e no target; os três componentes mais afetados \\
    \textbf{Fábrica}     & Ranking de ocorrências por linha
                         & Ocorrências do período, linha com mais ocorrências, linhas
                           monitoradas e período anterior; ranking das linhas de
                           montagem, com a variação; setores com mais ocorrências \\
    \textbf{Ofensores}   & Principais ofensores do período
                         & Os cinco defeitos e as oito linhas de maior IF Cost; o part
                           number e o defeito mais críticos \\
    \textbf{Prioridades} & Ocorrências prioritárias
                         & As ocorrências críticas de maior custo, com a descrição;
                           situação das revisões (pendentes, em revisão e
                           justificadas) e categorias predominantes \\
    \bottomrule
\end{tabularx}
\end{table}

\section{Encerramento da sessão}

O nome do usuário fica no pé do menu lateral. O botão \botao{\faEllipsisV}, ao lado
dele, leva ao \textbf{Perfil}. Ali, a seção \textbf{Sessões ativas} lista os
dispositivos conectados à conta: \botao{Encerrar sessão} desconecta um deles e
\botao{Encerrar outras sessões} desconecta todos, menos o atual.

Em computador compartilhado, encerre a sessão ao terminar e feche o navegador.

\fimdecapitulo


% ---------- CONSULTA DETALHADA ----------
% Os ajustes para blocos extensos começam somente aqui.
% Carregado somente depois do capítulo 2; preserva o layout do guia rápido.

% Extension for the complete user manual; original identity remains in the sty.
\makeatletter
\setcounter{tocdepth}{1}
% Explicit row shading avoids automatic row-color state across many longtables.
\renewcommand{\guiazebra}{}
% Chapter openings already delimit the content; omit a standalone closing badge.
\renewcommand{\fimdecapitulo}{}
\setlength{\guiarodapefolga}{0pt}
\setlength{\headheight}{34pt}
\setlength{\headsep}{10pt}
\let\guiaOldSubsection\subsection
\RenewDocumentCommand{\subsection}{s m}{%
 \Needspace{100pt}%
 \IfBooleanTF{#1}{\guiaOldSubsection*{#2}}{\guiaOldSubsection{#2}}}
\fancyhead[L]{\parbox[t]{0.49\textwidth}{\raggedright\guiacabecalhofonte\leftmark}}
\fancyhead[R]{\parbox[t]{0.47\textwidth}{\raggedleft\guiacabecalhofonte\rightmark}}
\titleformat{\part}[display]
 {\thispagestyle{empty}\vspace*{28pt}\raggedright}
 {{\guiahead\fontsize{130}{130}\selectfont\bfseries\color{guiavermelho!18}\thepart}}
 {-36pt}{\guiahead\fontsize{34}{40}\selectfont\bfseries\color{guiatinta}}
 [\vspace{12pt}{\color{guiavermelho}\rule{0.42\textwidth}{3bp}}\vspace{16pt}
 {\normalfont\fontsize{11}{16}\selectfont\color{guiacinza}\printcontents{l}{0}{\setcounter{tocdepth}{0}}}]
\newsavebox{\guiaPainelBox}
\NewDocumentCommand{\guiapainel}{O{} m m +m}{%
 \par\addvspace{7pt}%
 \begin{lrbox}{\guiaPainelBox}\begin{minipage}{\textwidth}
 \begin{minipage}[t]{\guiacolimagem}\vspace{0pt}
 \guiatela{#1}{#2}\par
 {\centering\fontsize{7.5}{10}\selectfont\color{guiacinza}
 \textbf{Figura \theguiapasso}\ifcsname guiacap@#2\endcsname
 ~— \csname guiacap@#2\endcsname\fi\par}
 \end{minipage}\hfill
 {\color{guiaborda}\vrule width 0.7pt}\hfill
 \begin{minipage}[t]{\guiacoltexto}\vspace{0pt}
 \setlength{\parskip}{0.35em}\raggedright\fontsize{9.3}{12.6}\selectfont
 {\guiahead\bfseries\fontsize{8}{8}\selectfont\colorbox{guiavermelho}{\color{white}\strut PASSO \theguiapasso}}\par
 \vspace{0.25em}
 {\guiahead\bfseries\fontsize{13.5}{16.5}\selectfont\color{guiatinta}#3\par}
 #4\par\end{minipage}\par\vspace{10pt}
 {\color{guiaborda}\hrule height 0.4pt}
 \end{minipage}\end{lrbox}
 \Needspace{\dimexpr\ht\guiaPainelBox+\dp\guiaPainelBox+16pt\relax}
 \ifguiaRotularPasso
  \phantomsection\label{passo:\theguiapasso}\guiaRotularPassofalse
 \fi
 \noindent\usebox{\guiaPainelBox}\par\vspace{10pt}
}
\newif\ifguiaRotularPasso
\RenewDocumentCommand{\passo}{O{} m m +m}{%
 \stepcounter{guiapasso}\protected@edef\@currentlabel{\theguiapasso}%
 \guiaRotularPassotrue
 \guiapainel[#1]{#2}{#3}{#4}}
\NewDocumentCommand{\passocontinuacao}{O{} m m +m}{%
 \guiapainel[#1]{#2}{#3}{#4}}

% Recortes e chamadas numeradas: aplicados somente depois dos capítulos originais.
\tikzset{guia marcador/.append style={minimum size=13pt,line width=0.9pt,
 font=\guiahead\bfseries\fontsize{7.5}{7.5}\selectfont}}
\newcommand{\guiaincluircaptura}[1]{%
 \ifcsname guiatela@trim@#1\endcsname
  \edef\guia@imageopts{width=\the\guiacolimagem,trim={\csname guiatela@trim@#1\endcsname},clip}%
  \expandafter\includegraphics\expandafter[\guia@imageopts]{#1}%
 \else\includegraphics[width=\guiacolimagem]{#1}\fi}
\renewcommand{\guiatela}[2]{%
 \begin{tikzpicture}
  \node[anchor=south west,inner sep=0pt] (tela) at (0,0) {\guiaincluircaptura{#2}};
  \draw[overlay,guiaborda,line width=0.7pt] (tela.south west) rectangle (tela.north east);
  \ifblank{#1}{}{%
   \begin{scope}[overlay,x={(tela.south east)},y={(tela.north west)}]
    \foreach \tx/\ty/\lx/\ly [count=\mn] in {#1}{%
     \draw[white,line width=1.8pt] (\lx,{1-\ly}) -- (\tx,{1-\ty});
     \draw[guiavermelho,line width=0.7pt] (\lx,{1-\ly}) -- (\tx,{1-\ty});
     \node[guia marcador] at (\lx,{1-\ly}) {\mn};
    }%
   \end{scope}%
  }%
 \end{tikzpicture}}

\expandafter\def\csname guiatela@trim@completo/02-perfil\endcsname{0.000bp 11.985bp 0.000bp 55.930bp}
\expandafter\def\csname guiatela@trim@completo/03-configuracoes\endcsname{0.000bp 12.090bp 0.000bp 56.420bp}
\expandafter\def\csname guiatela@trim@completo/04-classificacoes-material\endcsname{256.000bp 16.040bp 24.000bp 232.580bp}
\expandafter\def\csname guiatela@trim@completo/05-tipos-scrap\endcsname{264.000bp 12.075bp 24.000bp 241.500bp}
\expandafter\def\csname guiatela@trim@completo/06-metas-if-cost\endcsname{328.000bp 7.980bp 328.000bp 11.970bp}
\expandafter\def\csname guiatela@trim@completo/07-dados-producao\endcsname{460.800bp 43.945bp 464.000bp 47.940bp}
\expandafter\def\csname guiatela@trim@completo/08-programa-ano-identificacao\endcsname{256.000bp 1.626bp 16.000bp 231.705bp}
\expandafter\def\csname guiatela@trim@completo/09-programa-ano-execucao\endcsname{256.000bp 7.910bp 16.000bp 83.055bp}
\expandafter\def\csname guiatela@trim@completo/10-cotacao-dolar\endcsname{392.000bp 8.030bp 400.000bp 20.075bp}
\expandafter\def\csname guiatela@trim@completo/16-base-scrap-filtros\endcsname{864.000bp 169.470bp 56.000bp 209.820bp}
\expandafter\def\csname guiatela@trim@completo/17-ocorrencia-problema-local\endcsname{400.000bp 8.110bp 368.000bp 85.155bp}
\expandafter\def\csname guiatela@trim@completo/18-ocorrencia-causa-risco\endcsname{8.000bp 4.030bp 8.000bp 56.420bp}
\expandafter\def\csname guiatela@trim@completo/19-finalizar-analise\endcsname{592.000bp 283.605bp 600.000bp 303.030bp}
\expandafter\def\csname guiatela@trim@completo/21-salvar-modelo-favorito\endcsname{270.750bp 256.335bp 335.730bp 241.475bp}
\expandafter\def\csname guiatela@trim@completo/29-target-realizado\endcsname{288.000bp 24.210bp 272.000bp 56.490bp}
\expandafter\def\csname guiatela@trim@completo/30-dashboard-scrap-rate\endcsname{272.000bp 16.140bp 16.000bp 76.665bp}
\expandafter\def\csname guiatela@trim@completo/31-dashboard-quantidade\endcsname{272.000bp 11.880bp 16.000bp 75.240bp}
\expandafter\def\csname guiatela@trim@completo/32-dashboard-periodo\endcsname{1040.000bp 495.320bp 216.000bp 81.200bp}
\expandafter\def\csname guiatela@trim@completo/34-dashboard-filtros\endcsname{592.000bp 120.150bp 176.000bp 168.210bp}
\expandafter\def\csname guiatela@trim@completo/35-relatorios\endcsname{248.000bp 289.800bp 8.000bp 72.450bp}
\expandafter\def\csname guiatela@trim@completo/36-novo-relatorio-analises\endcsname{304.000bp 32.080bp 300.800bp 32.080bp}
\expandafter\def\csname guiatela@trim@completo/38-novo-relatorio-dados\endcsname{304.000bp 32.120bp 300.800bp 32.120bp}
\expandafter\def\csname guiatela@trim@completo/39-novo-relatorio-5-porques\endcsname{304.000bp 32.160bp 300.800bp 32.160bp}
\expandafter\def\csname guiatela@trim@completo/40-novo-relatorio-acoes\endcsname{304.000bp 32.240bp 300.800bp 32.240bp}
\expandafter\def\csname guiatela@trim@completo/41-relatorio-visao-geral\endcsname{248.000bp 88.880bp 24.000bp 129.280bp}
\expandafter\def\csname guiatela@trim@completo/42-relatorio-investigacao\endcsname{248.000bp 11.985bp 16.000bp 103.870bp}
\expandafter\def\csname guiatela@trim@completo/43-tarefas-lista\endcsname{248.000bp 96.600bp 16.000bp 120.750bp}
\expandafter\def\csname guiatela@trim@completo/44-tarefas-quadro\endcsname{248.000bp 12.135bp 16.000bp 137.530bp}
\expandafter\def\csname guiatela@trim@completo/45-tarefa-criacao-rapida\endcsname{4.743bp 6.760bp 1193.655bp 517.140bp}
\expandafter\def\csname guiatela@trim@completo/46-tarefa-validar-conclusao\endcsname{816.000bp 469.800bp 384.000bp 147.900bp}
\expandafter\def\csname guiatela@trim@completo/47-nova-tarefa-dados\endcsname{382.400bp 48.420bp 376.000bp 52.455bp}
\expandafter\def\csname guiatela@trim@completo/48-nova-tarefa-responsavel-prazo\endcsname{395.200bp 39.050bp 384.000bp 46.860bp}
\expandafter\def\csname guiatela@trim@completo/50-tarefas-cronograma\endcsname{272.000bp 11.430bp 16.000bp 121.920bp}
\expandafter\def\csname guiatela@trim@completo/51-relatorio-resultado\endcsname{288.000bp 12.105bp 16.000bp 121.050bp}
\expandafter\def\csname guiatela@trim@completo/52-relatorio-fotos\endcsname{272.000bp 11.970bp 16.000bp 87.780bp}
\expandafter\def\csname guiatela@trim@completo/53-relatorio-previa\endcsname{272.000bp 11.835bp 8.000bp 82.845bp}
\expandafter\def\csname guiatela@trim@completo/54-exportar-pdf\endcsname{352.000bp 24.180bp 364.800bp 20.150bp}
\expandafter\def\csname guiatela@trim@completo/55-exportar-powerpoint\endcsname{352.000bp 24.180bp 364.800bp 20.150bp}
\expandafter\def\csname guiatela@trim@completo/56-slide-exportado\endcsname{208.000bp 19.975bp 8.000bp 55.930bp}
\expandafter\def\csname guiatela@trim@completo/57-alertas\endcsname{264.000bp 11.985bp 8.000bp 79.900bp}
\expandafter\def\csname guiatela@trim@completo/58-configurar-alertas\endcsname{512.000bp 233.160bp 512.000bp 225.120bp}
\expandafter\def\csname guiatela@trim@completo/59-regra-alerta-condicoes\endcsname{432.000bp 20.175bp 432.000bp 16.140bp}
\expandafter\def\csname guiatela@trim@completo/60-regra-alerta-destinatarios\endcsname{376.000bp 8.100bp 424.000bp 16.200bp}
\expandafter\def\csname guiatela@trim@completo/61-integracao-teams\endcsname{120.000bp 48.300bp 8.000bp 64.400bp}
\expandafter\def\csname guiatela@trim@completo/62-teams-novo-destino-ingles\endcsname{120.000bp 12.045bp 16.000bp 56.210bp}
\expandafter\def\csname guiatela@trim@completo/65-fale-com-suporte\endcsname{516.800bp 144.900bp 518.400bp 140.875bp}
\expandafter\def\csname guiatela@trim@completo/66-solicitacoes-suporte\endcsname{128.000bp 12.225bp 16.000bp 81.500bp}
\expandafter\def\csname guiatela@trim@completo/67-usuarios\endcsname{120.000bp 32.320bp 8.000bp 64.640bp}
\expandafter\def\csname guiatela@trim@completo/68-novo-usuario\endcsname{512.000bp 72.810bp 512.000bp 72.810bp}
\expandafter\def\csname guiatela@trim@completo/69-uso-plataforma\endcsname{120.000bp 12.225bp 8.000bp 65.200bp}
\expandafter\def\csname guiatela@trim@completo/70-tempos-manuais\endcsname{512.000bp 143.820bp 512.000bp 151.810bp}
\makeatother
\expandafter\def\csname guiacap@completo/01-login\endcsname{Tela inicial de acesso ao sistema Hanaro}
\expandafter\def\csname guiacap@completo/02-perfil\endcsname{Tela de Perfil do colaborador e gerenciamento de sessões conectadas}
\expandafter\def\csname guiacap@completo/03-configuracoes\endcsname{Tela geral de Configurações com seleção de tema visual, idioma e regras do sistema}
\expandafter\def\csname guiacap@completo/04-classificacoes-material\endcsname{Tela de Classificações de Material com códigos pendentes e regras cadastradas}
\expandafter\def\csname guiacap@completo/05-tipos-scrap\endcsname{Tela de cadastro e gerenciamento de Tipos de Scrap}
\expandafter\def\csname guiacap@completo/06-metas-if-cost\endcsname{Janela de configuração da meta anual e mensal de IF Cost}
\expandafter\def\csname guiacap@completo/07-dados-producao\endcsname{Janela de lançamento mensal de Dados de produção}
\expandafter\def\csname guiacap@completo/08-programa-ano-identificacao\endcsname{Parte superior da tela Programa do ano com Identificação e KPI}
\expandafter\def\csname guiacap@completo/09-programa-ano-execucao\endcsname{Parte inferior da tela Programa do ano com Plano de execução e Organização}
\expandafter\def\csname guiacap@completo/10-cotacao-dolar\endcsname{Janela de consulta e ajuste da Cotação do dólar}
\expandafter\def\csname guiacap@completo/11-execucoes\endcsname{Tela de Execuções e atualização com o Histórico de rotinas processadas}
\expandafter\def\csname guiacap@completo/12-enviar-relatorio\endcsname{Janela para selecionar e processar manualmente o relatório GERP}
\expandafter\def\csname guiacap@completo/13-base-scrap\endcsname{Tela principal da Base de Scrap listando as ocorrências importadas}
\expandafter\def\csname guiacap@completo/14-base-scrap-ordenacao\endcsname{Opções de ordenação da tabela de Ocorrências de scrap}
\expandafter\def\csname guiacap@completo/15-base-scrap-status\endcsname{Filtro rápido por Status da análise na Base de Scrap}
\expandafter\def\csname guiacap@completo/16-base-scrap-filtros\endcsname{Painel de Filtros da base aberto na categoria Produção}
\expandafter\def\csname guiacap@completo/17-ocorrencia-problema-local\endcsname{Seções 1 (O problema) e 2 (Onde e quando) do formulário de Análise de Scrap}
\expandafter\def\csname guiacap@completo/18-ocorrencia-causa-risco\endcsname{Seções 3 (Causa) e 4 (Risco e responsabilidade) do formulário de Análise de Scrap}
\expandafter\def\csname guiacap@completo/19-finalizar-analise\endcsname{Janela de confirmação da finalização da análise}
\expandafter\def\csname guiacap@completo/20-ocorrencia-revisada\endcsname{Visualização resumida de uma ocorrência com análise finalizada (Revisado)}
\expandafter\def\csname guiacap@completo/21-salvar-modelo-favorito\endcsname{Janela Salvar como modelo favorito com definição do nome de referência}
\expandafter\def\csname guiacap@completo/22-nome-modelo-favorito\endcsname{Ajuste do nome de referência em Salvar como modelo favorito}
\expandafter\def\csname guiacap@completo/23-modelos-salvos\endcsname{Painel de Modelos Salvos aberto e modo de aplicação em lote ativado}
\expandafter\def\csname guiacap@completo/24-selecionar-ocorrencias\endcsname{Seleção de múltiplas ocorrências na tabela para aplicação do modelo}
\expandafter\def\csname guiacap@completo/25-aplicar-modelo-2-ocorrencias\endcsname{Confirmação de aplicação de modelo em lote para 2 ocorrências selecionadas}
\expandafter\def\csname guiacap@completo/26-aplicar-modelo-3-ocorrencias\endcsname{Confirmação de aplicação de modelo em lote para 3 ocorrências selecionadas}
\expandafter\def\csname guiacap@completo/27-modelo-lote-concluido\endcsname{Mensagem de conclusão após criar análises em lote com sucesso}
\expandafter\def\csname guiacap@completo/28-dashboard-financeiro\endcsname{Dashboard na visão financeira padrão (Impacto absoluto em Dólares)}
\expandafter\def\csname guiacap@completo/29-target-realizado\endcsname{Tabela detalhada Target \ensuremath{\times} realizado mensal com notas explicativas por mês}
\expandafter\def\csname guiacap@completo/30-dashboard-scrap-rate\endcsname{Dashboard com a visão de Eficiência relativa (Scrap Rate por mês e por produto)}
\expandafter\def\csname guiacap@completo/31-dashboard-quantidade\endcsname{Dashboard na visão de Quantidade física em unidades e acompanhamento semanal}
\expandafter\def\csname guiacap@completo/32-dashboard-periodo\endcsname{Menu de seleção do período de comparação histórica no topo do Dashboard}
\expandafter\def\csname guiacap@completo/33-dashboard-top-10\endcsname{Gráficos de ranking Top 10 linhas, modelos, componentes e ofensores}
\expandafter\def\csname guiacap@completo/34-dashboard-filtros\endcsname{Janela de Filtros do Dashboard aberta na categoria Operação}
\expandafter\def\csname guiacap@completo/35-relatorios\endcsname{Tela principal de Relatórios com abas de status e lista de planos}
\expandafter\def\csname guiacap@completo/36-novo-relatorio-analises\endcsname{Etapa 1 de criação de Novo relatório antes da seleção das análises}
\expandafter\def\csname guiacap@completo/37-novo-relatorio-analises-selecionadas\endcsname{Etapa 1 de criação de Novo relatório com análises selecionadas}
\expandafter\def\csname guiacap@completo/38-novo-relatorio-dados\endcsname{Etapa 2 de criação de Novo relatório (Título, Descrição, Responsável e Prazo)}
\expandafter\def\csname guiacap@completo/39-novo-relatorio-5-porques\endcsname{Etapa 3 de criação de Novo relatório com preenchimento dos 5 Porquês}
\expandafter\def\csname guiacap@completo/40-novo-relatorio-acoes\endcsname{Etapa 4 de criação de Novo relatório para inclusão opcional de ações iniciais}
\expandafter\def\csname guiacap@completo/41-relatorio-visao-geral\endcsname{Aba Visão geral do relatório com indicadores do caso e checklist de Prontidão}
\expandafter\def\csname guiacap@completo/42-relatorio-investigacao\endcsname{Aba Investigação do relatório com Causa raiz, 5 porquês e análises vinculadas}
\expandafter\def\csname guiacap@completo/43-tarefas-lista\endcsname{Aba Lista de tarefas do relatório com filtros e barra de adição rápida}
\expandafter\def\csname guiacap@completo/44-tarefas-quadro\endcsname{Aba Quadro de tarefas do relatório com as quatro colunas de andamento}
\expandafter\def\csname guiacap@completo/45-tarefa-criacao-rapida\endcsname{Criação rápida de tarefa diretamente na base da coluna Planejada}
\expandafter\def\csname guiacap@completo/46-tarefa-validar-conclusao\endcsname{Cartão de tarefa na coluna Em validação com botão Validar conclusão}
\expandafter\def\csname guiacap@completo/47-nova-tarefa-dados\endcsname{Parte superior da janela Nova tarefa (Dados da tarefa e Classificação)}
\expandafter\def\csname guiacap@completo/48-nova-tarefa-responsavel-prazo\endcsname{Parte inferior da janela Nova tarefa (Responsável, prazo, linhas e economia)}
\expandafter\def\csname guiacap@completo/49-tarefa-area-prazo\endcsname{Detalhe do preenchimento de Área responsável e Prazo na janela de tarefa}
\expandafter\def\csname guiacap@completo/50-tarefas-cronograma\endcsname{Aba Cronograma exibindo a linha do tempo de execução das tarefas}
\expandafter\def\csname guiacap@completo/51-relatorio-resultado\endcsname{Aba Resultado com comparativo antes e depois e medição de efeito na base de scrap}
\expandafter\def\csname guiacap@completo/52-relatorio-fotos\endcsname{Evidências visuais do relatório: fotos antes e depois e registros das etapas do processo}
\expandafter\def\csname guiacap@completo/53-relatorio-previa\endcsname{Aba Prévia do relatório com menu superior de ações aberto}
\expandafter\def\csname guiacap@completo/54-exportar-pdf\endcsname{Janela Montar documento configurada para exportar relatório executivo em PDF A4}
\expandafter\def\csname guiacap@completo/55-exportar-powerpoint\endcsname{Janela Montar documento configurada para exportar Apresentação 16:9 em PowerPoint com seleção de idioma}
\expandafter\def\csname guiacap@completo/56-slide-exportado\endcsname{Exemplo de slide exportado para PowerPoint com problema, causa raiz, ações e resultados}
\expandafter\def\csname guiacap@completo/57-alertas\endcsname{Tela principal de Alertas com o Histórico de notificações recebidas}
\expandafter\def\csname guiacap@completo/58-configurar-alertas\endcsname{Janela Configurar alertas listando as regras de notificação cadastradas}
\expandafter\def\csname guiacap@completo/59-regra-alerta-condicoes\endcsname{Parte superior do cadastro de Regra de notificação (Identificação e Quando avisar)}
\expandafter\def\csname guiacap@completo/60-regra-alerta-destinatarios\endcsname{Parte inferior do cadastro de Regra de notificação (Filtros, Quem recebe e Canais)}
\expandafter\def\csname guiacap@completo/61-integracao-teams\endcsname{Tela inicial de Integrações com o Microsoft Teams antes do cadastro de destinos}
\expandafter\def\csname guiacap@completo/62-teams-novo-destino-ingles\endcsname{Formulário de cadastro de novo destino do Microsoft Teams com idioma em Inglês}
\expandafter\def\csname guiacap@completo/63-teams-tipos-destino\endcsname{Opções do campo Tipo no cadastro de destino do Microsoft Teams (Pessoa, Grupo ou Canal)}
\expandafter\def\csname guiacap@completo/64-teams-novo-destino-portugues\endcsname{Formulário de cadastro de novo destino do Microsoft Teams com idioma em Português}
\expandafter\def\csname guiacap@completo/65-fale-com-suporte\endcsname{Janela Fale com o suporte para envio de dúvidas, problemas ou sugestões}
\expandafter\def\csname guiacap@completo/66-solicitacoes-suporte\endcsname{Tela administrativa de Solicitações de suporte (Caixa de entrada)}
\expandafter\def\csname guiacap@completo/67-usuarios\endcsname{Tela de Usuários listando os colaboradores cadastrados e seus cargos}
\expandafter\def\csname guiacap@completo/68-novo-usuario\endcsname{Janela Novo usuário preenchida para criação de conta de acesso}
\expandafter\def\csname guiacap@completo/69-uso-plataforma\endcsname{Painel de Uso da plataforma com tempo economizado e gráfico de acessos por dia}
\expandafter\def\csname guiacap@completo/70-tempos-manuais\endcsname{Janela Configurar tempos manuais em minutos por tarefa}
\expandafter\def\csname guiacap@completo/71-modo-tv\endcsname{Modo TV na tela Fábrica com indicadores, ranking e controles de rotação}

\expandafter\def\csname guiatela@trim@ia/01-assistente-lista\endcsname{0bp 55bp 0bp 0bp}
\expandafter\def\csname guiacap@ia/01-assistente-lista\endcsname{Esquema ilustrativo: Lista do Assistente}
\expandafter\def\csname guiatela@trim@ia/02-preanalise-leitura\endcsname{0bp 55bp 0bp 0bp}
\expandafter\def\csname guiacap@ia/02-preanalise-leitura\endcsname{Esquema ilustrativo: Leitura da pré-análise}
\expandafter\def\csname guiatela@trim@ia/03-preanalise-decisao\endcsname{0bp 55bp 0bp 0bp}
\expandafter\def\csname guiacap@ia/03-preanalise-decisao\endcsname{Esquema ilustrativo: Decisões sobre a pré-análise}
\expandafter\def\csname guiatela@trim@ia/04-consulta\endcsname{0bp 55bp 0bp 0bp}
\expandafter\def\csname guiacap@ia/04-consulta\endcsname{Esquema ilustrativo: Consulta a análises aprovadas}
\expandafter\def\csname guiatela@trim@ia/05-memoria\endcsname{0bp 55bp 0bp 0bp}
\expandafter\def\csname guiacap@ia/05-memoria\endcsname{Esquema ilustrativo: Memória do Assistente}

\part{Consulta Detalhada}
\chapter{Acesso ao Sistema e Gestão da Conta}

\label{cap:manual-1}

\section{Perfis de acesso}


Este guia apresenta o uso das principais telas do sistema \textbf{Hanaro}, orientando passo a passo as rotinas diárias de controle de refugo (\emph{scrap}), acompanhamento de metas financeiras (\emph{IF Cost}), investigação de causas, gestão de planos de ação e configuração de alertas.\par
As telas e menus visíveis para cada colaborador dependem do seu \textbf{Cargo (Perfil de Acesso)} cadastrado no sistema:\par
\begingroup
\fontsize{9}{12}\selectfont
\setlength{\LTpre}{4pt}
\setlength{\LTpost}{8pt}
\begin{longtable}{@{}L{2.425cm}L{4.244cm}L{8.488cm}@{}}
\toprule
\guiacabtabela \thead{Perfil de Acesso} & \thead{Quem Utiliza} & \thead{Telas Disponíveis no Menu Lateral} \\
\midrule
\endfirsthead
\toprule
\guiacabtabela \thead{Perfil de Acesso} & \thead{Quem Utiliza} & \thead{Telas Disponíveis no Menu Lateral} \\
\midrule
\endhead
\bottomrule
\endlastfoot
\rowcolor{white} \textbf{Analista} & Equipe operacional e técnica de Qualidade e Manufatura & \textbf{Dashboard}, \textbf{Execuções}, \textbf{Base de Scrap}, \textbf{Relatórios}, \textbf{Alertas}, \textbf{Configurações}, \textbf{Ajuda e suporte} e \textbf{Perfil} \\
\rowcolor{guiaareia!70} \textbf{Gestor} & Líderes de área, coordenadores e gerentes industriais & As rotinas do Analista e a aprovação de tarefas em validação, além dos recursos de gestão liberados para sua conta \\
\rowcolor{white} \textbf{Administrador (Admin)} & Responsáveis pela administração do sistema e integrações & Área administrativa: \textbf{Uso da plataforma}, \textbf{Execuções}, \textbf{Suporte} e \textbf{Perfil}. \textbf{Usuários} e outros cadastros dependem das permissões administrativas da conta \\
\end{longtable}
\endgroup


\begin{remark}[info]
Analista e Gestor atuam nas rotinas da planta. A conta Admin pertence à área de
administração da plataforma e não representa automaticamente acesso a todas as telas
operacionais. A gestão de usuários exige permissão de superusuário. Assistente,
Consulta, Memória e integrações podem variar conforme o perfil e o ambiente.
\end{remark}

\clearpage
\section{Tela de Login}

\subsection*{Objetivo}

Permitir que o colaborador entre no sistema Hanaro com seu usuário e senha individuais. Deve ser utilizada por todos os colaboradores no início do expediente ou quando a sessão for encerrada por inatividade.\par

\subsection*{Antes de começar}

\begin{itemize}
\item Possuir usuário e senha ativos cadastrados pela administração do sistema.
\end{itemize}

\begin{remark}[info]
O procedimento de acesso já está descrito no guia rápido. As tabelas abaixo complementam as regras de credenciais e as mensagens de erro. Consulte o passo~\ref{passo:1} (\pagina{passo:1}).
\end{remark}

\Needspace{160pt}

\subsection*{Campos e regras}

\begingroup
\fontsize{9}{12}\selectfont
\setlength{\LTpre}{4pt}
\setlength{\LTpost}{8pt}
\begin{longtable}{@{}L{3.389cm}L{1.916cm}L{2.652cm}L{6.778cm}@{}}
\toprule
\guiacabtabela \thead{Campo} & \thead{Obrigatório} & \thead{Formato} & \thead{Descrição e regra de preenchimento} \\
\midrule
\endfirsthead
\toprule
\guiacabtabela \thead{Campo} & \thead{Obrigatório} & \thead{Formato} & \thead{Descrição e regra de preenchimento} \\
\midrule
\endhead
\bottomrule
\endlastfoot
\rowcolor{white} \textbf{Usuário} & Sim & Texto (letras minúsculas e números) & Identificação única do colaborador no sistema. \\
\rowcolor{guiaareia!70} \textbf{Senha} & Sim & Texto oculto por pontos & Senha pessoal de acesso. Possui botão lateral com ícone de olho para exibir ou ocultar o texto digitado. \\
\end{longtable}
\endgroup

\subsection*{Mensagens e solução de problemas}

\begin{remark}
\textbf{Informe o nome de usuário e a senha.}\par
\textbf{Causa provável:} Você clicou em \textbf{Entrar} deixando o campo \textbf{Usuário} ou \textbf{Senha} em branco.\par
\textbf{Como resolver:} Preencha os dois campos antes de clicar em \textbf{Entrar}.
\end{remark}

\begin{remark}
\textbf{Nome de usuário ou senha incorretos.}\par
\textbf{Causa provável:} O usuário não existe, está desativado ou a senha foi digitada com algum erro.\par
\textbf{Como resolver:} Clique no ícone de olho para conferir a senha digitada. Se continuar sem acesso, solicite a redefinição de senha ao Administrador.
\end{remark}

\begin{remark}
\textbf{Caps Lock ativado.}\par
\textbf{Causa provável:} A tecla de letras maiúsculas (\emph{Caps Lock}) do teclado está ativada enquanto você digita a senha.\par
\textbf{Como resolver:} Desative a tecla \emph{Caps Lock} do teclado e digite a senha novamente.
\end{remark}

\begin{remark}
\textbf{Muitas tentativas...}\par
\textbf{Causa provável:} Foram feitas várias tentativas seguidas com senha incorreta.\par
\textbf{Como resolver:} Aguarde alguns minutos pelo desbloqueio automático antes de tentar novamente.
\end{remark}

\Needspace{100pt}
\begin{remark}
\textbf{Não foi possível entrar. Verifique a conexão e tente novamente.}\par
\textbf{Causa provável:} O computador está sem comunicação com a rede da empresa ou o sistema está temporariamente indisponível.\par
\textbf{Como resolver:} Verifique sua conexão de rede e tente acessar novamente em instantes.
\end{remark}

\section{Meu Perfil e Sessões Ativas}

\subsection*{Objetivo}

Permitir que qualquer colaborador consulte seus dados cadastrais, atualize sua foto de identificação, ative o recebimento de avisos pessoais no Microsoft Teams e encerre acessos abertos em outros computadores.\par

\subsection*{Antes de começar}

\begin{itemize}
\item Estar conectado ao sistema.
\item Para trocar a foto de identificação, possuir uma imagem salva no computador nos formatos \tech{JPEG}, \tech{PNG} ou \tech{WebP} com tamanho máximo de \textbf{5 MB}.
\end{itemize}

\passo[0.06700/0.95847/0.07900/0.90929,0.39000/0.13115/0.39000/0.13115,0.18200/0.40874/0.24000/0.36066,0.87500/0.55628/0.84300/0.49727,0.94100/0.55628/0.96500/0.61749,0.94500/0.71913/0.91000/0.66120,0.93200/0.86776/0.87500/0.83060]{completo/02-perfil}{Consultar perfil e sessões}{%
\begin{elementos}
\item \textbf{Nome e cargo.} Clique no seu nome, na parte inferior do menu lateral, para abrir o Perfil.
\item \textbf{Identificação.} Confira nome completo, usuário, e-mail e cargo cadastrados.
\item \textbf{Alertas no Teams.} Ative ou desative Receber meus alertas no Teams conforme a disponibilidade do seu destino pessoal.
\item \textbf{Alterar foto.} Escolha um arquivo JPEG, PNG ou WebP de até 5 MB. Aguarde a mensagem Foto atualizada com sucesso.
\item \textbf{Remover foto.} Apague a foto atual se desejar retirar a identificação visual.
\item \textbf{Atualizar sessões.} Confira os computadores e navegadores conectados e recarregue a lista. A indicação Este dispositivo identifica o computador em uso.
\item \textbf{Encerrar sessão.} Desconecte imediatamente o acesso selecionado.
\end{elementos}
}

\Needspace{160pt}

\subsection*{Campos e regras}

\begingroup
\fontsize{9}{12}\selectfont
\setlength{\LTpre}{4pt}
\setlength{\LTpost}{8pt}
\begin{longtable}{@{}L{3.389cm}L{1.916cm}L{2.652cm}L{6.778cm}@{}}
\toprule
\guiacabtabela \thead{Campo} & \thead{Obrigatório} & \thead{Formato} & \thead{Descrição e regra de preenchimento} \\
\midrule
\endfirsthead
\toprule
\guiacabtabela \thead{Campo} & \thead{Obrigatório} & \thead{Formato} & \thead{Descrição e regra de preenchimento} \\
\midrule
\endhead
\bottomrule
\endlastfoot
\rowcolor{white} \textbf{Nome completo} & Somente leitura & Texto & Nome completo do colaborador cadastrado pela administração. \\
\rowcolor{guiaareia!70} \textbf{Usuário} & Somente leitura & Texto (identificação de acesso) & Identificação utilizada pelo colaborador para entrar no sistema. \\
\rowcolor{white} \textbf{E-mail da conta} & Somente leitura & Endereço de e-mail & E-mail corporativo vinculado à conta do colaborador. \\
\rowcolor{guiaareia!70} \textbf{Receber meus alertas no Teams} & Não & Chave liga/desliga & Ativa o envio de alertas individuais para o seu Microsoft Teams quando houver um destino pessoal cadastrado. \\
\rowcolor{white} \textbf{Cargo} & Somente leitura & Texto (\tech{Analista}, \tech{Gestor} ou \tech{Admin}) & Indica o nível de permissão do colaborador no sistema. \\
\rowcolor{guiaareia!70} \textbf{Foto de identificação} & Não & Imagem (\tech{JPEG}, \tech{PNG} ou \tech{WebP} até 5 MB) & Foto exibida no menu lateral e nos registros de atividades do colaborador. \\
\rowcolor{white} \textbf{Sessões ativas} & Somente leitura & Lista de dispositivos & Mostra navegador, sistema operacional, data/hora do último acesso e endereço de rede de cada sessão aberta. \\
\end{longtable}
\endgroup

\subsection*{Mensagens e solução de problemas}

\begin{remark}[info]
\textbf{Foto atualizada com sucesso.}\par
\textbf{O que significa:} A nova foto escolhida foi salva corretamente.\par
\textbf{Próximo passo:} Nenhuma ação necessária (mensagem de confirmação).
\end{remark}

\begin{remark}
\textbf{Use uma imagem JPEG, PNG ou WebP de até 5 MB.}\par
\textbf{Causa provável:} O arquivo escolhido está em outro formato (como PDF ou BMP) ou pesa mais de 5 MB.\par
\textbf{Como resolver:} Escolha uma imagem \tech{.jpg} ou \tech{.png} menor que 5 MB e clique novamente em \textbf{Alterar foto}.
\end{remark}

\begin{remark}
\textbf{Nenhum destino pessoal configurado}\par
\textbf{Causa provável:} A chave do Teams está visível, mas ainda não há um canal pessoal vinculado ao seu usuário.\par
\textbf{Como resolver:} Solicite ao Administrador o cadastro do seu destino individual na tela de \textbf{Integrações > Microsoft Teams}.
\end{remark}

\fimdecapitulo

\chapter{Configurações e Parametrizações Operacionais}

\label{cap:manual-2}

\section{Painel Geral de Configurações (Tema Visual e Idioma)}

\subsection*{Objetivo}

Permitir ao usuário escolher entre o tema visual claro ou escuro, alterar o idioma de exibição do sistema (Português, Inglês ou Coreano) e acessar os seis painéis de regras e cadastros operacionais da fábrica.\par

\subsection*{Antes de começar}

\begin{itemize}
\item Estar conectado ao sistema e acessar o menu lateral \textbf{Configurações}.
\end{itemize}

\passo[0.07000/0.31475/0.07000/0.31475,0.36900/0.24044/0.36900/0.24044,0.54300/0.42186/0.51000/0.36393,0.54300/0.52131/0.51000/0.56831,0.78000/0.36066/0.78000/0.36066,0.57500/0.74317/0.57500/0.74317]{completo/03-configuracoes}{Ajustar tema e idioma}{%
\begin{elementos}
\item \textbf{Configurações.} Abra o menu de engrenagem.
\item \textbf{Claro e Escuro.} Escolha o cartão de tema adequado ao ambiente.
\item \textbf{Tema escuro.} Use a chave para alternar diretamente entre claro e escuro.
\item \textbf{Usar configuração do sistema.} Ative para acompanhar o tema do computador.
\item \textbf{Idioma do sistema.} Escolha Português (BR), English ou 한국어. A captura ilustra a seleção de coreano; os controles permanecem nas posições indicadas.
\item \textbf{Dados e regras.} Os cartões abrem Classificações de Material, Tipos de Scrap, IF Cost, Dados de produção, Programa do ano e Cotação do dólar, detalhados nos passos seguintes.
\end{elementos}
}

\Needspace{160pt}

\subsection*{Campos e regras}

\begingroup
\fontsize{9}{12}\selectfont
\setlength{\LTpre}{4pt}
\setlength{\LTpost}{8pt}
\begin{longtable}{@{}L{3.389cm}L{1.916cm}L{2.652cm}L{6.778cm}@{}}
\toprule
\guiacabtabela \thead{Campo} & \thead{Obrigatório} & \thead{Formato} & \thead{Descrição e regra de preenchimento} \\
\midrule
\endfirsthead
\toprule
\guiacabtabela \thead{Campo} & \thead{Obrigatório} & \thead{Formato} & \thead{Descrição e regra de preenchimento} \\
\midrule
\endhead
\bottomrule
\endlastfoot
\rowcolor{white} \textbf{Tema da interface (``Claro'' / ``Escuro'')} & Sim & Seleção visual / Chave liga/desliga & Altera imediatamente as cores de fundo e textos de todas as telas. \\
\rowcolor{guiaareia!70} \textbf{Usar configuração do sistema} & Não & Chave liga/desliga & Faz o Hanaro seguir automaticamente o tema claro ou escuro configurado no Windows/Mac. \\
\rowcolor{white} \textbf{Idioma do sistema} & Sim & Seleção única (\tech{Português (BR)}, \tech{English}, {\color{guiavermelho}한국어}) & Define o idioma em que todos os botões, menus e títulos são exibidos para você. \\
\end{longtable}
\endgroup

\subsection*{Mensagens e solução de problemas}

\begin{remark}
\textbf{Menus e textos do sistema aparecem em coreano (``한국어'') ou Inglês}\par
\textbf{Causa provável:} Outro idioma foi selecionado acidentalmente no quadro superior direito de Configurações.\par
\textbf{Como resolver:} Clique no último item com ícone de engrenagem no menu lateral esquerdo e selecione a primeira opção do quadro superior direito (\textbf{Português (BR)}).
\end{remark}

\section{Classificações de Material}

\subsection*{Objetivo}

Traduzir os códigos contábeis e de itens importados do sistema GERP em nomes claros e definir quais contas entram ou ficam fora do cálculo de \emph{IF Cost} no Dashboard, na Base de Scrap e nos Relatórios. Deve ser operada por \textbf{Gestores} e \textbf{Administradores}.\par

\subsection*{Antes de começar}

\begin{itemize}
\item Ter realizado ao menos uma importação de relatório do GERP para que o sistema identifique eventuais códigos novos sem regra.
\end{itemize}

\passo[0.42424/0.22464/0.42424/0.22464,0.94182/0.39855/0.89697/0.46377,0.93939/0.08696/0.88485/0.03623,0.29818/0.67246/0.29091/0.73913,0.94424/0.56522/0.99030/0.47826,0.83394/0.56522/0.77576/0.49275,0.93576/0.84783/0.84242/0.78986,0.96000/0.84783/0.95152/0.96377]{completo/04-classificacoes-material}{Classificar materiais}{%
\begin{elementos}
\item \textbf{Pendências e filtros.} Em Códigos sem regra, filtre Tipo, Situação, Falta e Registros ou busque um código importado do GERP.
\item \textbf{Criar regra.} Na linha pendente, cadastre o nome ou a classificação que falta.
\item \textbf{Verificar de novo.} Atualize a lista de códigos sem regra.
\item \textbf{Buscar regras.} Use o Tipo ou a busca por código ou resultado no quadro Regras.
\item \textbf{Nova regra.} Cadastre uma tradução manualmente.
\item \textbf{Reaplicar ao histórico.} Atualize os registros antigos. Uma regra nova já vale para as próximas importações, mas não altera o histórico automaticamente.
\item \textbf{Editar.} Use o lápis da linha para alterar a regra existente.
\item \textbf{Excluir.} Use a lixeira da linha para remover a regra.
\end{elementos}
}

\Needspace{160pt}

\subsection*{Campos e regras}

\begingroup
\fontsize{9}{12}\selectfont
\setlength{\LTpre}{4pt}
\setlength{\LTpost}{8pt}
\begin{longtable}{@{}L{3.389cm}L{1.916cm}L{2.652cm}L{6.778cm}@{}}
\toprule
\guiacabtabela \thead{Campo} & \thead{Obrigatório} & \thead{Formato} & \thead{Descrição e regra de preenchimento} \\
\midrule
\endfirsthead
\toprule
\guiacabtabela \thead{Campo} & \thead{Obrigatório} & \thead{Formato} & \thead{Descrição e regra de preenchimento} \\
\midrule
\endhead
\bottomrule
\endlastfoot
\rowcolor{white} \textbf{Filtros de Códigos sem regra (``Tipo'', ``Situação'', ``Falta'', ``Registros'')} & Não & Campos de seleção & Ajudam a localizar códigos vindos do GERP que ainda estão sem nome do item ou sem conta classificada. \\
\rowcolor{guiaareia!70} \textbf{Tipo (da Regra)} & Sim & Campo de seleção (ex.: \tech{Conta do IF Cost}, \tech{Nome do item}) & Indica qual informação do arquivo GERP será traduzida pela regra. \\
\rowcolor{white} \textbf{Código de origem} & Sim & Texto exato do GERP & Código original vindo do relatório GERP (por exemplo, \tech{DEFECT - DIRECT CHARGED} ou \tech{CONSUMABLE SUPPLIES EXPENSE}). \\
\rowcolor{guiaareia!70} \textbf{Resultado} & Sim & Texto / Classificação & Define o nome amigável do item ou se a conta \textbf{Entra no IF Cost} / fica \textbf{Fora do IF Cost}. \\
\end{longtable}
\endgroup

\subsection*{Mensagens e solução de problemas}

\begin{remark}
\textbf{Códigos antigos continuam sem classificação no Dashboard após criar uma regra nova}\par
\textbf{Causa provável:} A regra foi salva, mas ainda não foi aplicada aos lotes importados anteriormente.\par
\textbf{Como resolver:} Clique no botão \textbf{Reaplicar ao histórico} no topo da lista de \textbf{Regras}.
\end{remark}

\section{Tipos de Scrap}

\subsection*{Objetivo}

Manter a lista padronizada de categorias de defeito ou perda (como \tech{Montagem}, \tech{Componente}, \tech{Processo}, \tech{Oxidação}, \tech{Avaria de transporte}) que os analistas escolhem ao analisar uma ocorrência na Base de Scrap.\par

\subsection*{Antes de começar}

\begin{itemize}
\item Possuir acesso ao menu \textbf{Configurações}.
\end{itemize}

\passo[0.11585/0.18978/0.06707/0.12409,0.30122/0.18978/0.34756/0.12409,0.92317/0.35766/0.87195/0.27737,0.65122/0.69635/0.59146/0.62774,0.94268/0.69635/0.85366/0.80292,0.96829/0.69635/0.95366/0.94161]{completo/05-tipos-scrap}{Gerenciar tipos de scrap}{%
\begin{elementos}
\item \textbf{Nome.} Informe a categoria de scrap, como Oxidação ou Avaria de transporte.
\item \textbf{Código.} Confira o valor gerado a partir do nome. Use letras maiúsculas, números, \_ ou -. Ajuste antes de adicionar: o código não poderá ser alterado depois.
\item \textbf{Adicionar tipo.} Grave a nova categoria.
\item \textbf{Disponível nas análises.} Libere ou oculte o tipo nos formulários de novas análises.
\item \textbf{Editar.} Renomeie a categoria com o lápis.
\item \textbf{Excluir.} Remova a categoria com a lixeira.
\end{elementos}
}

\Needspace{160pt}

\subsection*{Campos e regras}

\begingroup
\fontsize{9}{12}\selectfont
\setlength{\LTpre}{4pt}
\setlength{\LTpost}{8pt}
\begin{longtable}{@{}L{3.389cm}L{1.916cm}L{2.652cm}L{6.778cm}@{}}
\toprule
\guiacabtabela \thead{Campo} & \thead{Obrigatório} & \thead{Formato} & \thead{Descrição e regra de preenchimento} \\
\midrule
\endfirsthead
\toprule
\guiacabtabela \thead{Campo} & \thead{Obrigatório} & \thead{Formato} & \thead{Descrição e regra de preenchimento} \\
\midrule
\endhead
\bottomrule
\endlastfoot
\rowcolor{white} \textbf{Nome *} & Sim & Texto livre & Nome da categoria exibido para toda a equipe durante a análise de ocorrências. \\
\rowcolor{guiaareia!70} \textbf{Código *} & Sim & Letras maiúsculas, números, \tech{\_} ou \tech{-} & Identificador fixo da categoria (ex.: \tech{MONTAGEM}, \tech{PROCESSO}). Não pode ser alterado após a criação. \\
\rowcolor{white} \textbf{Disponível nas análises} & Sim & Chave liga/desliga & Controla se a categoria aparece como opção para novos registros na Base de Scrap. \\
\end{longtable}
\endgroup

\subsection*{Mensagens e solução de problemas}

\begin{remark}
\textbf{Campo Código bloqueado durante a edição de um tipo existente}\par
\textbf{Causa provável:} Por segurança do histórico de relatórios já emitidos, o código identificador é fixo após a criação.\par
\textbf{Como resolver:} Se precisar de um código diferente, desative a chave \textbf{Disponível nas análises} do item antigo e crie um novo tipo.
\end{remark}

\section{Metas Anuais e Mensais de IF Cost}

\subsection*{Objetivo}

Definir as metas mensais de \emph{IF Cost} em dólares (\tech{US\$}) para o ano, permitindo distribuir os valores automaticamente pelos 12 meses para comparação direta com o custo de scrap no Dashboard.\par

\subsection*{Antes de começar}

\begin{itemize}
\item Contas contábeis do \emph{IF Cost} já classificadas em \textbf{Configurações > Classificações de Material}.
\end{itemize}

\passo[0.11017/0.20923/0.09831/0.16923,0.48305/0.38462/0.48305/0.28205,0.86271/0.28410/0.93051/0.23590,0.24576/0.56410/0.27966/0.49744,0.16102/0.66359/0.09322/0.71282,0.39492/0.66359/0.39831/0.71282,0.59831/0.66359/0.68644/0.63077,0.14407/0.79897/0.09322/0.87692]{completo/06-metas-if-cost}{Definir metas de IF Cost}{
Em Configurações, abra IF Cost.\par%
\begin{elementos}
\item \textbf{Ano.} Selecione o ano de planejamento.
\item \textbf{Meta do ano.} Confira Total planejado, Média mensal e a variação em relação ao ano anterior.
\item \textbf{Ver contas do IF Cost.} Consulte as contas que entram na composição da meta.
\item \textbf{Modo de preenchimento.} Escolha Reduzir de janeiro a dezembro ou Mesmo valor todo mês.
\item \textbf{Meta de janeiro.} Informe o valor inicial, por exemplo US\$ 12.000,00.
\item \textbf{Meta de dezembro.} No modo de redução gradual, informe o valor final, por exemplo US\$ 6.000,00.
\item \textbf{Preencher tabela.} Distribua os valores pelos doze meses.
\item \textbf{Cartões mensais.} Confira os valores e ajuste os meses necessários.
\end{elementos}
\par\smallskip {\fontsize{8.2}{11}\selectfont\color{guiacinza}Role até o fim da janela e clique em Salvar metas para gravar os valores.\par}
}

\Needspace{160pt}

\subsection*{Campos e regras}

\begingroup
\fontsize{9}{12}\selectfont
\setlength{\LTpre}{4pt}
\setlength{\LTpost}{8pt}
\begin{longtable}{@{}L{3.389cm}L{1.916cm}L{2.652cm}L{6.778cm}@{}}
\toprule
\guiacabtabela \thead{Campo} & \thead{Obrigatório} & \thead{Formato} & \thead{Descrição e regra de preenchimento} \\
\midrule
\endfirsthead
\toprule
\guiacabtabela \thead{Campo} & \thead{Obrigatório} & \thead{Formato} & \thead{Descrição e regra de preenchimento} \\
\midrule
\endhead
\bottomrule
\endlastfoot
\rowcolor{white} \textbf{Ano} & Sim & Ano (\tech{AAAA}) & Define o ano das metas de \emph{IF Cost}. \\
\rowcolor{guiaareia!70} \textbf{Modo de preenchimento} & Não & Botões de opção & Escolhe entre redução progressiva de janeiro a dezembro ou valor fixo mensal. \\
\rowcolor{white} \textbf{Meta de janeiro / Meta de dezembro} & Não & Valor em dólares (\tech{US\$}) & Valores inicial e final usados pelo botão \textbf{\ensuremath{\checkmark} Preencher tabela}. \\
\rowcolor{guiaareia!70} \textbf{Valores de ``Jan.'' a ``Dez.''} & Sim & Valor em dólares (\tech{US\$}) & Meta financeira de \emph{IF Cost} de cada mês do ano. \\
\end{longtable}
\endgroup

\subsection*{Mensagens e solução de problemas}

\begin{remark}
\textbf{Os valores calculados para os 12 meses não ficaram gravados ao fechar a janela}\par
\textbf{Causa provável:} O botão \textbf{\ensuremath{\checkmark} Preencher tabela} apenas preenche os campos na tela; é necessário confirmar a gravação no final da janela.\par
\textbf{Como resolver:} Após clicar em \textbf{\ensuremath{\checkmark} Preencher tabela}, role até o rodapé da janela e clique em \textbf{Salvar metas}.
\end{remark}

\section{Dados de Produção Mensal}

\subsection*{Objetivo}

Registrar o valor financeiro produzido (em dólares) e a quantidade de peças fabricadas em cada mês. O Dashboard e os Relatórios utilizam esses números para calcular a taxa de refugo em relação à produção (\emph{Scrap Rate}).\par

\subsection*{Antes de começar}

\begin{itemize}
\item Ter em mãos os números de fechamento mensal de produção da fábrica.
\end{itemize}

\passo[0.25355/0.23616/0.12322/0.19209,0.58294/0.23616/0.64455/0.18079,0.30806/0.31525/0.38389/0.33898,0.27962/0.52768/0.12322/0.48588,0.72986/0.52768/0.88152/0.48588,0.50000/0.63503/0.54976/0.68927,0.48578/0.87571/0.47867/0.81356,0.82701/0.93333/0.88152/0.98305,0.58531/0.93333/0.52607/0.98305]{completo/07-dados-producao}{Registrar dados de produção}{
Em Configurações, abra Dados de produção.\par%
\begin{elementos}
\item \textbf{Escopo.} Selecione Global (todos os produtos) ou um produto específico.
\item \textbf{Ano.} Selecione o ano de lançamento.
\item \textbf{Mês.} Selecione o mês a preencher. Os meses preenchidos exibem um ponto verde.
\item \textbf{Valor da produção.} Informe o total produzido em dólares.
\item \textbf{Quantidade produzida.} Informe as unidades fabricadas no mês.
\item \textbf{Observação.} Registre, se necessário, uma nota do fechamento. Use Enter para avançar entre os campos de valor, quantidade e o mês seguinte.
\item \textbf{Resumo acumulado.} Confira Valor no ano, Quantidade no ano e Detalhado por produto.
\item \textbf{Salvar dados.} Grave os lançamentos do ano indicado no botão.
\item \textbf{Limpar dados.} Use apenas quando precisar apagar os lançamentos do ano selecionado.
\end{elementos}
}

\Needspace{160pt}

\subsection*{Campos e regras}

\begingroup
\fontsize{9}{12}\selectfont
\setlength{\LTpre}{4pt}
\setlength{\LTpost}{8pt}
\begin{longtable}{@{}L{3.389cm}L{1.916cm}L{2.652cm}L{6.778cm}@{}}
\toprule
\guiacabtabela \thead{Campo} & \thead{Obrigatório} & \thead{Formato} & \thead{Descrição e regra de preenchimento} \\
\midrule
\endfirsthead
\toprule
\guiacabtabela \thead{Campo} & \thead{Obrigatório} & \thead{Formato} & \thead{Descrição e regra de preenchimento} \\
\midrule
\endhead
\bottomrule
\endlastfoot
\rowcolor{white} \textbf{Escopo} & Sim & Campo de seleção & Define se o lançamento é do total da fábrica (\tech{Global}) ou por linha de produto. \\
\rowcolor{guiaareia!70} \textbf{Ano} & Sim & Ano (\tech{AAAA}) & Ano de referência da produção. \\
\rowcolor{white} \textbf{Valor da produção (USD)} & Sim & Valor em dólares (\tech{US\$}) & Base utilizada pelo sistema para calcular o percentual financeiro de \emph{Scrap Rate}. \\
\rowcolor{guiaareia!70} \textbf{Quantidade produzida} & Sim & Número inteiro & Total de peças produzidas no mês selecionado. \\
\rowcolor{white} \textbf{Observação} & Não & Texto livre & Comentário opcional sobre eventos produtivos daquele mês. \\
\end{longtable}
\endgroup

\subsection*{Mensagens e solução de problemas}

\begin{remark}
\textbf{Diferença entre "Valor no ano" e "Detalhado por produto"}\par
\textbf{Causa provável:} A soma dos valores lançados por produto está diferente do valor lançado no escopo Global.\par
\textbf{Como resolver:} Confira o indicador percentual abaixo do campo de observação e ajuste os valores para que fechem em 100\% antes de salvar.
\end{remark}

\section{Programa do Ano}

\subsection*{Objetivo}

Preencher a ficha anual do programa de redução de scrap da fábrica. As informações salvas nesta tela compõem automaticamente a abertura executiva dos relatórios exportados no ano.\par

\subsection*{Antes de começar}

\begin{itemize}
\item Metas de \emph{IF Cost} do ano já salvas (para que o sistema apresente a referência automática do indicador).
\end{itemize}

\passo[0.06506/0.06592/0.03012/0.03506,0.46988/0.28892/0.46988/0.28892,0.10241/0.42216/0.06024/0.37167,0.30120/0.42216/0.31325/0.35764,0.49398/0.42216/0.55422/0.37167,0.48193/0.84432/0.74699/0.77840,0.04819/0.91304/0.10843/0.86255,0.22169/0.91304/0.28916/0.86255,0.42771/0.91304/0.54217/0.89762,0.45422/0.95652/0.74699/0.94670]{completo/08-programa-ano-identificacao}{Identificar o programa do ano}{
Em Configurações, abra Programa do ano.\par%
\begin{elementos}
\item \textbf{Ano.} Escolha o ano do programa e confira a indicação de ficha salva.
\item \textbf{Tema.} Informe o objetivo anual, como reduzir o IF Cost em 15\% sobre o ano anterior.
\item \textbf{Período.} Informe a vigência, por exemplo janeiro a dezembro.
\item \textbf{Equipe.} Identifique as áreas participantes.
\item \textbf{Local.} Informe a planta ou o local do programa.
\item \textbf{Indicador.} O KPI usa os dados da base e das metas. Substitua o indicador automático somente se necessário.
\item \textbf{Nível atual.} Role para preencher o campo completo quando for usar um valor manual.
\item \textbf{Meta.} Role para conferir ou substituir a meta automática.
\item \textbf{Nota da meta.} Registre a explicação da meta, quando aplicável.
\item \textbf{Voltar ao automático.} Restaure os textos e valores obtidos da base e das metas.
\end{elementos}
\par\smallskip {\fontsize{8.2}{11}\selectfont\color{guiacinza}Os marcadores inferiores apontam para os rótulos visíveis do KPI. Role a página para mostrar os campos completos antes de preenchê-los.\par}
}

\passo[0.12048/0.18644/0.10843/0.12429,0.32651/0.18644/0.32530/0.12429,0.65060/0.18644/0.68675/0.12429,0.06627/0.46328/0.26506/0.48588,0.96627/0.18644/0.98313/0.28814,0.10120/0.66667/0.07229/0.61582,0.30120/0.66667/0.28916/0.61582,0.49639/0.66667/0.51807/0.61582,0.24819/0.81356/0.22892/0.75706,0.74096/0.81356/0.74699/0.75706,0.93133/0.93785/0.85542/0.93220]{completo/09-programa-ano-execucao}{Planejar a execução do programa}{%
\begin{elementos}
\item \textbf{Etapa.} Nomeie cada etapa do plano. O programa aceita até seis etapas.
\item \textbf{Período.} Informe quando a etapa será executada.
\item \textbf{O que é feito.} Descreva o trabalho previsto para a etapa.
\item \textbf{Adicionar etapa.} Inclua outra linha no plano.
\item \textbf{Excluir etapa.} Remova a linha com a lixeira correspondente.
\item \textbf{Patrocinador.} Informe quem patrocina o programa.
\item \textbf{Responsável.} Informe o responsável pelo programa.
\item \textbf{Líder.} Identifique quem coordena a execução.
\item \textbf{Tempo integral.} Informe os participantes dedicados integralmente.
\item \textbf{Tempo parcial.} Informe os demais participantes.
\item \textbf{Salvar ficha.} Grave o programa do ano.
\end{elementos}
}

\Needspace{160pt}

\subsection*{Campos e regras}

\begingroup
\fontsize{9}{12}\selectfont
\setlength{\LTpre}{4pt}
\setlength{\LTpost}{8pt}
\begin{longtable}{@{}L{3.389cm}L{1.916cm}L{2.652cm}L{6.778cm}@{}}
\toprule
\guiacabtabela \thead{Campo} & \thead{Obrigatório} & \thead{Formato} & \thead{Descrição e regra de preenchimento} \\
\midrule
\endfirsthead
\toprule
\guiacabtabela \thead{Campo} & \thead{Obrigatório} & \thead{Formato} & \thead{Descrição e regra de preenchimento} \\
\midrule
\endhead
\bottomrule
\endlastfoot
\rowcolor{white} \textbf{Tema / Período / Equipe / Local} & Não & Texto livre & Identificação geral do programa anual. Campos deixados em branco não aparecem no relatório impresso. \\
\rowcolor{guiaareia!70} \textbf{Indicador / Nível atual / Meta / Nota da meta} & Não & Texto livre & Em branco, usa os valores automáticos das metas do ano; preencha apenas se quiser substituir manualmente. \\
\rowcolor{white} \textbf{Plano de execução (``Etapa'', ``Período'', ``O que é feito'')} & Não & Lista de até 6 etapas & Descreve as fases do cronograma anual de redução de scrap. \\
\rowcolor{guiaareia!70} \textbf{Organização (``Patrocinador'', ``Responsável'', ``Líder'', ``Tempo integral'', ``Tempo parcial'')} & Não & Texto livre & Registra a liderança e os integrantes dedicados ao programa. \\
\end{longtable}
\endgroup

\subsection*{Mensagens e solução de problemas}

\begin{remark}
\textbf{Não é possível adicionar mais etapas no Plano de execução}\par
\textbf{Causa provável:} O limite máximo de 6 etapas da ficha anual já foi atingido.\par
\textbf{Como resolver:} Agrupe atividades semelhantes dentro das 6 etapas disponíveis.
\end{remark}

\section{Cotação do Dólar}

\subsection*{Objetivo}

Consultar a cotação diária de venda do dólar (PTAX do Banco Central) usada para converter valores nos envios manuais de relatórios e permitir ajustes manuais de taxa quando necessário.\par

\subsection*{Antes de começar}

\begin{itemize}
\item Possuir acesso ao menu \textbf{Configurações}.
\end{itemize}

\passo[0.11881/0.23316/0.08911/0.18653,0.41584/0.51813/0.41584/0.51813,0.91089/0.37824/0.95050/0.33161]{completo/10-cotacao-dolar}{Consultar a cotação do dólar}{
Em Configurações, abra Cotação do dólar.\par%
\begin{elementos}
\item \textbf{Ano.} Use as setas para escolher o período.
\item \textbf{Cotações.} Confira Data, R\$ por US\$ com quatro casas decimais e Origem. Automação identifica a atualização automática.
\item \textbf{Editar.} Use o lápis da data desejada para informar e salvar uma taxa manual. O valor manual tem prioridade sobre a cotação automática.
\end{elementos}
\par\smallskip {\fontsize{8.2}{11}\selectfont\color{guiacinza}A captura mostra a lista, antes de abrir a edição da linha.\par}
}

\Needspace{160pt}

\subsection*{Campos e regras}

\begingroup
\fontsize{9}{12}\selectfont
\setlength{\LTpre}{4pt}
\setlength{\LTpost}{8pt}
\begin{longtable}{@{}L{3.389cm}L{1.916cm}L{2.652cm}L{6.778cm}@{}}
\toprule
\guiacabtabela \thead{Campo} & \thead{Obrigatório} & \thead{Formato} & \thead{Descrição e regra de preenchimento} \\
\midrule
\endfirsthead
\toprule
\guiacabtabela \thead{Campo} & \thead{Obrigatório} & \thead{Formato} & \thead{Descrição e regra de preenchimento} \\
\midrule
\endhead
\bottomrule
\endlastfoot
\rowcolor{white} \textbf{Data} & Somente leitura & Data (\tech{AAAA-MM-DD}) & Data do lançamento. Em sábados, domingos e feriados, o sistema utiliza automaticamente a cotação do dia útil anterior. \\
\rowcolor{guiaareia!70} \textbf{R\$ por US\$} & Sim (ao editar) & Número decimal (4 casas, ex.: \tech{5.1500}) & Taxa de conversão de 1 dólar para reais. \\
\rowcolor{white} \textbf{Origem} & Somente leitura & Texto (\tech{Automação} / Manual) & Mostra se a cotação veio da atualização automática do Banco Central ou se foi digitada por um usuário. \\
\end{longtable}
\endgroup

\subsection*{Mensagens e solução de problemas}

\begin{remark}
\textbf{Lançamento de sábado ou domingo usou a cotação de sexta-feira}\par
\textbf{Causa provável:} Em fins de semana e feriados não há boletim PTAX, por isso o sistema aplica automaticamente a taxa do último dia útil anterior.\par
\textbf{Como resolver:} Comportamento padrão do sistema. Se precisar de outra taxa para aquela data, clique no ícone de lápis da linha e digite o valor manualmente.
\end{remark}

\fimdecapitulo

\chapter{Execuções e Atualização de Dados GERP}

\label{cap:manual-3}

\section{Histórico de Rotinas e Cargas}

\subsection*{Objetivo}

Acompanhar o histórico de todas as cargas de dados importadas do sistema GERP, conferindo horários, tempo de processamento, quantidade de linhas recebidas, válidas ou rejeitadas e a confirmação de que os dados já estão publicados no sistema.\par

\subsection*{Antes de começar}

\begin{itemize}
\item Possuir perfil \textbf{Gestor} ou \textbf{Administrador} com acesso ao menu \textbf{Execuções}.
\end{itemize}

\begin{remark}[info]
A abertura do histórico, a conferência do status e o uso dos filtros já estão descritos no guia rápido. Consulte o passo~\ref{passo:3} (\pagina{passo:3}).
\end{remark}

\Needspace{160pt}

\subsection*{Referência das colunas}

A descrição das colunas está na \autoref{tab:historico}.

\subsection*{Mensagens e solução de problemas}

\begin{remark}
\textbf{Coluna REJEITADOS com número maior que zero (``> 0'')}\par
\textbf{Causa provável:} Algumas linhas do arquivo exportado do GERP estavam incompletas ou com formato inválido.\par
\textbf{Como resolver:} Clique sobre a linha da rotina para verificar o motivo da rejeição, corrija a extração no GERP e envie o arquivo novamente.
\end{remark}

\section{Envio Manual de Relatório GERP}

\subsection*{Objetivo}

Permitir o envio manual do arquivo \tech{Other Account Transaction Text} exportado do GERP para atualizar imediatamente a Base de Scrap e os gráficos do Dashboard.\par

\subsection*{Antes de começar}

\begin{itemize}
\item Ter exportado do GERP o relatório \tech{Other Account Transaction Text} no formato \textbf{TSV} (texto separado por tabulação, codificação \tech{UTF-8} ou \tech{CP1252}).
\end{itemize}

\begin{remark}[info]
A seleção e o envio do arquivo GERP já estão descritos no guia rápido. Use esta seção para conferir requisitos e mensagens de validação. Consulte o passo~\ref{passo:2} (\pagina{passo:2}).
\end{remark}

\Needspace{160pt}

\subsection*{Campos e regras}

\begingroup
\fontsize{9}{12}\selectfont
\setlength{\LTpre}{4pt}
\setlength{\LTpost}{8pt}
\begin{longtable}{@{}L{3.389cm}L{1.916cm}L{2.652cm}L{6.778cm}@{}}
\toprule
\guiacabtabela \thead{Campo} & \thead{Obrigatório} & \thead{Formato} & \thead{Descrição e regra de preenchimento} \\
\midrule
\endfirsthead
\toprule
\guiacabtabela \thead{Campo} & \thead{Obrigatório} & \thead{Formato} & \thead{Descrição e regra de preenchimento} \\
\midrule
\endhead
\bottomrule
\endlastfoot
\rowcolor{white} \textbf{Selecionar relatório GERP} & Sim & Arquivo \tech{.tsv} (\tech{UTF-8} ou \tech{CP1252}) & Arquivo original \tech{Other Account Transaction Text} exportado do GERP. \\
\end{longtable}
\endgroup

\subsection*{Mensagens e solução de problemas}

\begin{remark}
\textbf{Arquivo recusado ou falha ao processar relatório}\par
\textbf{Causa provável:} Foi selecionado um arquivo em formato incorreto (como \tech{.xlsx} ou \tech{.pdf}) ou uma planilha com colunas diferentes do padrão do GERP.\par
\textbf{Como resolver:} Exporte novamente o relatório \tech{Other Account Transaction Text} do GERP em formato \tech{.tsv} e repita o envio.
\end{remark}

\fimdecapitulo

\chapter{Base de Scrap e Análise de Ocorrências}

\label{cap:manual-4}

\section{Consulta, Ordenação e Filtros da Base de Scrap}

\subsection*{Objetivo}

Consultar todas as ocorrências de refugo importadas para o sistema, permitindo ordenar por data ou valor, filtrar por status de análise e localizar rapidamente os itens que precisam de investigação técnica.\par

\subsection*{Antes de começar}

\begin{itemize}
\item Estar conectado ao sistema e acessar o menu lateral \textbf{Base de Scrap}.
\end{itemize}

\begin{remark}[info]
A abertura da Base de Scrap e seus controles de ordenação e status já estão descritos no guia rápido. Consulte o passo~\ref{passo:5} (\pagina{passo:5}).
\end{remark}
\subsection*{Opções de consulta}
\begin{itemize}
\item A ordenação aceita \textbf{Data da transação}, \textbf{Organização}, \textbf{Código do item}, \textbf{Quantidade} e \textbf{Valor (USD)}, em ordem crescente ou decrescente.
\item O status \textbf{Não analisado} identifica ocorrências aguardando análise; \textbf{Rascunho}, análises iniciadas e salvas parcialmente; \textbf{Revisado}, análises concluídas.
\end{itemize}


\passo[0.12000/0.37547/0.04706/0.28302,0.44941/0.25849/0.37647/0.16981,0.79529/0.25849/0.87059/0.16981,0.44941/0.43774/0.37647/0.53774,0.79529/0.43774/0.87059/0.53774,0.44941/0.61887/0.37647/0.70755,0.92000/0.92453/0.91765/0.83019,0.79765/0.92453/0.75294/0.83019]{completo/16-base-scrap-filtros}{Aplicar filtros da base}{%
Abra Filtros da base pelo botão Filtrar. O detalhe mostra a categoria Produção.\par
\begin{elementos}
\item \textbf{Categorias.} Alterne Geral, Catálogo, Produção e Análise.
\item \textbf{Produção.} Selecione o produto ou o grupo de produção.
\item \textbf{Departamento.} Restrinja o recorte ao departamento desejado.
\item \textbf{Divisão.} Escolha a divisão, quando necessário.
\item \textbf{Filtro da linha.} Defina o agrupamento de linhas.
\item \textbf{Linha.} Escolha a linha de produção.
\item \textbf{Aplicar.} Atualize a tabela com o recorte selecionado.
\item \textbf{Limpar filtros.} Remova as restrições e volte à lista completa.
\end{elementos}
}

\Needspace{160pt}

\subsection*{Campos e regras}

\begingroup
\fontsize{9}{12}\selectfont
\setlength{\LTpre}{4pt}
\setlength{\LTpost}{8pt}
\begin{longtable}{@{}L{3.389cm}L{1.916cm}L{2.652cm}L{6.778cm}@{}}
\toprule
\guiacabtabela \thead{Campo} & \thead{Obrigatório} & \thead{Formato} & \thead{Descrição e regra de preenchimento} \\
\midrule
\endfirsthead
\toprule
\guiacabtabela \thead{Campo} & \thead{Obrigatório} & \thead{Formato} & \thead{Descrição e regra de preenchimento} \\
\midrule
\endhead
\bottomrule
\endlastfoot
\rowcolor{white} \textbf{Ordenar por} & Não & Campo de seleção & Organiza a tabela por \tech{Data da transação}, \tech{Organização}, \tech{Código do item}, \tech{Quantidade} ou \tech{Valor (USD)}. \\
\rowcolor{guiaareia!70} \textbf{Status da análise} & Não & Campo de seleção & Filtra entre \tech{Todos}, \tech{Não analisado}, \tech{Rascunho} e \tech{Revisado}. \\
\rowcolor{white} \textbf{Ordem (``Decrescente'' / ``Crescente'')} & Não & Botão alternador & Inverte a listagem do maior para o menor ou do mais recente para o mais antigo. \\
\rowcolor{guiaareia!70} \textbf{Filtros da base (``Geral'', ``Catálogo'', ``Produção'', ``Análise'')} & Não & Painel de filtros & Permite restringir a consulta por linha de produção, departamento, item ou analista responsável. \\
\end{longtable}
\endgroup

\subsection*{Mensagens e solução de problemas}

\begin{remark}
\textbf{Algumas ocorrências não aparecem na tabela}\par
\textbf{Causa provável:} Há um filtro ativo no painel \textbf{Filtrar} ou o seletor de status não está marcado em \textbf{Todos}.\par
\textbf{Como resolver:} Clique em \textbf{Filtrar > Limpar filtros} e selecione \textbf{Todos} no filtro de status da análise.
\end{remark}

\section{Preenchimento e Finalização de Análise de Scrap Individual}

\subsection*{Objetivo}

Registrar a investigação técnica de uma ocorrência de scrap em 7 etapas (O problema, Onde e quando, Causa, Risco e responsabilidade, Custo, Contenção imediata e Evidências/fotos), permitindo salvar como rascunho ou finalizar a análise.\par

\subsection*{Antes de começar}

\begin{itemize}
\item Categorias de scrap já cadastradas em \textbf{Configurações > Tipos de Scrap}.
\end{itemize}

\passo[0.46154/0.26554/0.21154/0.23164,0.46154/0.40113/0.21154/0.46893,0.15385/0.65537/0.09615/0.61582,0.38462/0.65537/0.40385/0.60452,0.61538/0.65537/0.67308/0.60452,0.24038/0.75367/0.11538/0.77401,0.65385/0.75367/0.75000/0.77401]{completo/17-ocorrencia-problema-local}{Descrever o problema e o local}{%
Na Base de Scrap, abra uma ocorrência Não analisada ou em Rascunho para preencher a análise.\par
\begin{elementos}
\item \textbf{Título da análise.} Resuma o defeito ou a causa.
\item \textbf{Descrição e causa raiz.} Explique o ocorrido, as causas identificadas e as ações tomadas. O campo aceita até 20.000 caracteres.
\item \textbf{Data.} Informe a data da ocorrência.
\item \textbf{Hora.} Informe o horário.
\item \textbf{Turno.} Selecione o turno.
\item \textbf{Linha.} Informe a linha.
\item \textbf{Configuração da linha.} Complete a configuração correspondente.
\end{elementos}
}

\passo[0.30808/0.19459/0.11616/0.15135,0.50000/0.29730/0.50000/0.23784,0.29798/0.41622/0.11616/0.45405,0.82323/0.37297/0.89394/0.32432,0.50000/0.62162/0.50000/0.56216,0.24747/0.74595/0.14646/0.70270,0.81313/0.74595/0.89394/0.70270,0.87374/0.96973/0.80303/0.90811,0.96667/0.96973/0.95455/0.90811]{completo/18-ocorrencia-causa-risco}{Registrar causa e responsabilidade}{%
\begin{elementos}
\item \textbf{Tipo de scrap.} Escolha a categoria da ocorrência.
\item \textbf{Categoria 4M.} Selecione Mão de obra, Máquina, Material, Método ou Outro.
\item \textbf{Causa analisada.} Descreva a causa investigada.
\item \textbf{Situação da causa.} Informe se é hipótese ou causa confirmada.
\item \textbf{Nível de risco.} Escolha Não avaliado, Baixo, Moderado, Alto ou Crítico.
\item \textbf{Onde/como.} Registre a segmentação do problema, por exemplo Empilhamento na saída da linha.
\item \textbf{Área responsável.} Informe a área responsável pela ocorrência.
\item \textbf{Salvar rascunho.} Guarde o preenchimento e continue depois.
\item \textbf{Finalizar relatório.} Conclua a análise quando todos os dados e evidências estiverem prontos.
\end{elementos}
\par\smallskip {\fontsize{8.2}{11}\selectfont\color{guiacinza}Antes de finalizar, role até Custo, Contenção imediata e Evidências e fotos. Confira Tipo de custo, Modelo e Quantidade. Se Quantidade ficar em branco, será usada a quantidade original da base. Registre a contenção e anexe até oito imagens JPEG, PNG ou WebP, com até 10 MB por foto.\par}
}

\passo[0.76863/0.75102/0.86275/0.73469]{completo/19-finalizar-analise}{Finalizar a análise}{
\begin{elementos}
\item \textbf{Finalizar relatório.} Na janela Finalizar análise, confira o aviso e clique em Finalizar relatório para confirmar. Após a finalização, a análise continuará editável apenas por seu autor.
\end{elementos}
}

\begin{remark}[info]
A consulta de uma ocorrência revisada já está descrita no guia rápido. Consulte o passo~\ref{passo:6} (\pagina{passo:6}).
\end{remark}
Depois de finalizar a análise, a etiqueta \textbf{Revisado} confirma a conclusão.
O ícone de lápis permite editar a análise; o coração abre o cadastro de modelo favorito,
detalhado na seção seguinte.


\Needspace{160pt}

\subsection*{Campos e regras}

\begingroup
\fontsize{9}{12}\selectfont
\setlength{\LTpre}{4pt}
\setlength{\LTpost}{8pt}
\begin{longtable}{@{}L{3.389cm}L{1.916cm}L{2.652cm}L{6.778cm}@{}}
\toprule
\guiacabtabela \thead{Campo} & \thead{Obrigatório} & \thead{Formato} & \thead{Descrição e regra de preenchimento} \\
\midrule
\endfirsthead
\toprule
\guiacabtabela \thead{Campo} & \thead{Obrigatório} & \thead{Formato} & \thead{Descrição e regra de preenchimento} \\
\midrule
\endhead
\bottomrule
\endlastfoot
\rowcolor{white} \textbf{Título da análise *} & Sim & Texto curto & Resumo breve que identifica a análise nas listas e relatórios. \\
\rowcolor{guiaareia!70} \textbf{Descrição e causa raiz *} & Sim & Texto longo (até 20.000 caracteres) & Detalhamento completo da ocorrência, causas e ações imediatas. \\
\rowcolor{white} \textbf{Data / Hora / Turno / Linha} & Não & Data (\tech{AAAA/MM/DD}), Hora e Seletores & Identificam o momento e o local da linha onde ocorreu a perda. \\
\rowcolor{guiaareia!70} \textbf{Tipo de scrap *} & Sim & Campo de seleção & Categoria escolhida dentre os \textbf{Tipos de Scrap} ativos. \\
\rowcolor{white} \textbf{Categoria 4M da causa} & Não & Botões de opção (\tech{Mão de obra}, \tech{Máquina}, \tech{Material}, \tech{Método}, \tech{Outro}) & Classifica a natureza da causa para análise nos relatórios. \\
\rowcolor{guiaareia!70} \textbf{Situação da causa / Nível de risco} & Não & Seletor e Botões de opção & Indicam se a causa ainda é uma hipótese e qual a gravidade do risco (\tech{Baixo} a \tech{Crítico}). \\
\rowcolor{white} \textbf{Onde/como (segmentação) / Área responsável} & Não & Texto livre & Detalham o posto ou etapa da linha e o setor responsável pela tratativa. \\
\rowcolor{guiaareia!70} \textbf{Tipo de custo (IF Cost) / Quantidade} & Não & Seletor e Número & Classificam o custo; quando a quantidade fica vazia, o sistema adota a quantidade da base. \\
\rowcolor{white} \textbf{Contenção imediata} & Não & Texto livre & Ação imediata executada no turno para interromper a perda. \\
\rowcolor{guiaareia!70} \textbf{Evidências e fotos (``0/8'')} & Não & Envio de imagens (\tech{JPEG}, \tech{PNG}, \tech{WebP}) & Permite anexar até \textbf{8 fotos} de no máximo \textbf{10 MB} cada. \\
\end{longtable}
\endgroup

\subsection*{Mensagens e solução de problemas}

\begin{remark}
\textbf{\textbullet{} Alterações não salvas (no canto inferior esquerdo)}\par
\textbf{Causa provável:} Você alterou informações na tela, mas ainda não clicou em \textbf{Salvar rascunho} nem em \textbf{Finalizar relatório}.\par
\textbf{Como resolver:} Clique em \textbf{Salvar rascunho} ou \textbf{Finalizar relatório} antes de fechar a janela.
\end{remark}

\begin{remark}
\textbf{O relatório será finalizado e continuará editável apenas por quem o criou. Deseja continuar?}\par
\textbf{Causa provável:} Aviso de segurança informando que, após finalizar, apenas o autor da análise poderá editá-la.\par
\textbf{Como resolver:} Clique em \textbf{Finalizar relatório} para confirmar. Caso outro colega ainda precise complementar dados antes, cancele e use \textbf{Salvar rascunho}.
\end{remark}

\begin{remark}
\textbf{Imagem recusada em Evidências e fotos}\par
\textbf{Causa provável:} A foto ultrapassa 10 MB, não está em \tech{JPEG}/\tech{PNG}/\tech{WebP} ou já existem 8 fotos anexadas.\par
\textbf{Como resolver:} Reduza o tamanho da imagem ou remova uma das 8 fotos antes de enviar outra.
\end{remark}

\section{Modelos Favoritos e Aplicação de Análise em Lote}

\subsection*{Objetivo}

Permitir que o analista salve uma análise recorrente como \textbf{Modelo Favorito} e aplique esse mesmo preenchimento de uma só vez a várias ocorrências semelhantes da Base de Scrap.\par

\subsection*{Antes de começar}

\begin{itemize}
\item Ter ao menos uma ocorrência já revisada (para salvar um modelo novo) ou um modelo já salvo na lista de \textbf{Modelos Salvos}.
\end{itemize}

\passo[0.50000/0.52121/0.20455/0.43939,0.83182/0.75758/0.88636/0.86364]{completo/21-salvar-modelo-favorito}{Salvar uma análise como modelo favorito}{
Abra uma ocorrência Revisada e clique no ícone de coração para abrir Salvar como modelo favorito.\par
\begin{elementos}
\item \textbf{Nome de referência.} Digite um nome que permita identificar o modelo na lista de Modelos Salvos.
\item \textbf{Salvar.} Grave o modelo para reutilizar o preenchimento em ocorrências semelhantes.
\end{elementos}
}

\begin{remark}[info]
A escolha do modelo, a seleção das ocorrências e a confirmação da aplicação em lote já estão descritas no guia rápido. Consulte o passo~\ref{passo:7} (\pagina{passo:7}).
\end{remark}
\subsection*{Validações e manutenção dos modelos}
\begin{itemize}
\item O painel \textbf{Modelos Salvos} também permite \textbf{Editar modelo} ou \textbf{Excluir} um modelo existente.
\item Antes da confirmação em lote, confira a quantidade de itens selecionados, o \textbf{Tipo de scrap} e a \textbf{Descrição e causa raiz} que serão aplicados.
\item A mensagem de criação dos relatórios confirma o término. Use \textbf{Fechar} para encerrar a janela.
\end{itemize}


\Needspace{160pt}

\subsection*{Campos e regras}

\begingroup
\fontsize{9}{12}\selectfont
\setlength{\LTpre}{4pt}
\setlength{\LTpost}{8pt}
\begin{longtable}{@{}L{3.389cm}L{1.916cm}L{2.652cm}L{6.778cm}@{}}
\toprule
\guiacabtabela \thead{Campo} & \thead{Obrigatório} & \thead{Formato} & \thead{Descrição e regra de preenchimento} \\
\midrule
\endfirsthead
\toprule
\guiacabtabela \thead{Campo} & \thead{Obrigatório} & \thead{Formato} & \thead{Descrição e regra de preenchimento} \\
\midrule
\endhead
\bottomrule
\endlastfoot
\rowcolor{white} \textbf{Nome de referência para este modelo} & Sim & Texto curto & Nome que identifica o modelo dentro da lista \textbf{Modelos Salvos}. \\
\rowcolor{guiaareia!70} \textbf{Caixas de marcação da tabela} & Sim (no modo lote) & Seleção múltipla & Marcam quais linhas da Base de Scrap receberão o preenchimento automático do modelo. \\
\end{longtable}
\endgroup

\subsection*{Mensagens e solução de problemas}

\begin{remark}
\textbf{Ocorrências inativas ou que já possuem relatório serão ignoradas.}\par
\textbf{Causa provável:} Entre os itens marcados na tabela, existe algum registro que já possui relatório vinculado.\par
\textbf{Como resolver:} Mensagem preventiva normal: o sistema aplicará o modelo apenas às ocorrências pendentes e preservará intactas as que já têm relatório.
\end{remark}

\begin{remark}[info]
\textbf{\ensuremath{\checkmark} 3 Relatórios criados com sucesso:}\par
\textbf{O que significa:} O modelo foi aplicado com sucesso a todas as ocorrências válidas selecionadas.\par
\textbf{Próximo passo:} Clique em \textbf{Fechar} para voltar à tabela atualizada.
\end{remark}

\fimdecapitulo

\chapter{Dashboard de Monitoramento e Indicadores}

\label{cap:manual-5}

\section{Visão financeira e tabela Target \ensuremath{\times} realizado}

\subsection*{Objetivo}

Acompanhar o resultado financeiro do refugo (\emph{IF Cost} em dólares) comparado às metas mensais e anuais da planta, permitindo abrir a tabela comparativa mês a mês com as notas explicativas de fechamento.\par

\subsection*{Antes de começar}

\begin{itemize}
\item Relatórios GERP importados em \textbf{Execuções} e metas anuais cadastradas em \textbf{Configurações > IF Cost}.
\end{itemize}

\begin{remark}[info]
A abertura do Dashboard e seus controles financeiros já estão descritos no guia rápido. Consulte o passo~\ref{passo:4} (\pagina{passo:4}).
\end{remark}
Os exemplos numéricos desta seção ilustram os dados das capturas e variam conforme o período e os filtros selecionados. YTD significa acumulado no ano.

\Needspace{190pt}
\subsection*{Interpretação dos cartões financeiros}
\begin{itemize}

\item \textbf{Acumulado:} Total de \emph{IF Cost} realizado no período (ex.: \tech{US\$ 94.176,73});
\item \textbf{Target de IF Cost:} Meta acumulada até o último mês com dados (ex.: \tech{US\$ 109.020,00});
\item \textbf{Atingimento do target:} Relação entre a meta e o custo realizado (meta $\div$ realizado $\times 100$). Quando o realizado é menor que a meta, o resultado supera 100\% (ex.: \tech{115,8\%}, favorável);
\item \textbf{Diferença para a meta:} Situação em relação à meta e folga em dólares (ex.: \tech{Dentro da meta \ensuremath{\checkmark} — Folga de US\$ 14.843,27});
\item \textbf{Variação de IF Cost:} Comparação percentual com o mesmo período do ano anterior (ex.: \tech{-24,7\% \ensuremath{\downarrow}}).

\end{itemize}

\passo[0.49231/0.54444/0.49231/0.54444,0.43077/0.57333/0.53846/0.62222,0.10154/0.30111/0.07692/0.25556,0.88462/0.30111/0.90769/0.25556]{completo/29-target-realizado}{Abrir a tabela Target \ensuremath{\times} realizado}{%
No gráfico do Dashboard, escolha Mensal ou Semanal e use o ícone de tabela para abrir este quadro.\par
\begin{elementos}
\item \textbf{Dados mensais.} Consulte Realizado, Referência, Meta, Atingimento e Variação por mês.
\item \textbf{Nota do mês.} Leia a observação que explica a variação do período.
\item \textbf{Tabela anterior.} Volte ao quadro anterior.
\item \textbf{Próxima tabela.} Avance entre as seis tabelas do Dashboard.
\end{elementos}
}

\Needspace{160pt}

\subsection*{Campos e controles da tela}

\begingroup
\fontsize{9}{12}\selectfont
\setlength{\LTpre}{4pt}
\setlength{\LTpost}{8pt}
\begin{longtable}{@{}L{3.389cm}L{1.916cm}L{2.652cm}L{6.778cm}@{}}
\toprule
\guiacabtabela \thead{Controle} & \thead{Obrigatório} & \thead{Formato} & \thead{Descrição e funcionamento} \\
\midrule
\endfirsthead
\toprule
\guiacabtabela \thead{Controle} & \thead{Obrigatório} & \thead{Formato} & \thead{Descrição e funcionamento} \\
\midrule
\endhead
\bottomrule
\endlastfoot
\rowcolor{white} \textbf{Chave ``Impacto absoluto'' / ``Eficiência relativa''} & Não & Chave de alternância & Alterna entre valores totais (\tech{US\$} ou \tech{Unidades}) e percentual sobre a produção (\tech{Scrap Rate \%}). \\
\rowcolor{guiaareia!70} \textbf{Chave ``Valor monetário'' / ``Quantidade''} & Não & Chave de alternância & Alterna a unidade de todos os gráficos entre Dólares (\tech{US\$}) e peças (\tech{Unidades}). \\
\rowcolor{white} \textbf{Botões ``Mensal'' / ``Semanal''} & Não & Botões de opção & Mudam a divisão do gráfico temporal entre meses (\tech{Jan–Dez}) e semanas (\tech{W01–W52}). \\
\rowcolor{guiaareia!70} \textbf{COMPARAR COM (na tabela)} & Não & Campo de seleção & Escolhe a base de comparação (\tech{YTD anterior}, \tech{Mesmo mês ano anterior} ou \tech{Mês anterior}). \\
\end{longtable}
\endgroup

\subsection*{Mensagens e solução de problemas}

\begin{remark}
\textbf{Meta (``Target'') aparece zerada nos cartões e no gráfico}\par
\textbf{Causa provável:} As metas mensais do ano selecionado ainda não foram salvas em \textbf{Configurações > IF Cost}.\par
\textbf{Como resolver:} Acesse \textbf{Configurações > IF Cost}, preencha as metas dos 12 meses e clique em \textbf{Salvar metas}.
\end{remark}

\section{Visão de Eficiência Relativa (Scrap Rate)}

\subsection*{Objetivo}

Analisar o custo de scrap proporcionalmente ao volume produzido em cada mês (\textbf{Scrap Rate de IF Cost}), permitindo comparar meses de diferentes volumes produtivos na mesma base percentual.\par

\subsection*{Antes de começar}

\begin{itemize}
\item Valores mensais de produção cadastrados em \textbf{Configurações > Dados de produção}.
\end{itemize}

\passo[0.05366/0.07458/0.03415/0.01695,0.14268/0.20000/0.09756/0.14124,0.39024/0.20000/0.35366/0.14124,0.63537/0.20000/0.62195/0.14124,0.87439/0.20000/0.89024/0.14124,0.32927/0.58192/0.32927/0.58192,0.86098/0.58192/0.86098/0.58192]{completo/30-dashboard-scrap-rate}{Consultar o Scrap Rate}{%
\begin{elementos}
\item \textbf{Eficiência relativa.} Ative a chave para consultar as taxas.
\item \textbf{Acumulado.} Leia o Scrap Rate de IF Cost, a perda percentual em relação à produção.
\item \textbf{IF Cost do recorte.} Confira o valor de scrap considerado no cálculo.
\item \textbf{Material Amount.} Confira o valor da produção do mesmo período.
\item \textbf{Variação da taxa.} Compare com o mesmo acumulado do ano anterior.
\item \textbf{Evolução mensal.} Acompanhe o gráfico Scrap Rate de IF Cost por mês.
\item \textbf{Produtos.} Compare as linhas de produto no ranking de Scrap Rate.
\end{elementos}
}

\Needspace{160pt}

\subsection*{Campos e indicadores da tela}

\begingroup
\fontsize{9}{12}\selectfont
\setlength{\LTpre}{4pt}
\setlength{\LTpost}{8pt}
\begin{longtable}{@{}L{3.389cm}L{1.916cm}L{2.652cm}L{6.778cm}@{}}
\toprule
\guiacabtabela \thead{Indicador} & \thead{Obrigatório} & \thead{Formato} & \thead{Descrição e funcionamento} \\
\midrule
\endfirsthead
\toprule
\guiacabtabela \thead{Indicador} & \thead{Obrigatório} & \thead{Formato} & \thead{Descrição e funcionamento} \\
\midrule
\endhead
\bottomrule
\endlastfoot
\rowcolor{white} \textbf{Scrap Rate de IF Cost} & Somente leitura & Percentual (\tech{0,0000\%}) & Calculado dividindo o \emph{IF Cost do recorte} pelo \emph{Material Amount} (valor produzido) e multiplicando o resultado por 100. \\
\rowcolor{guiaareia!70} \textbf{Material Amount} & Somente leitura & Valor em dólares (\tech{US\$}) & Soma do valor de produção cadastrado para os meses e produtos filtrados. \\
\end{longtable}
\endgroup

\subsection*{Mensagens e solução de problemas}

\begin{remark}
\textbf{Taxa percentual não calculada para determinado mês}\par
\textbf{Causa provável:} O mês consultado ainda não possui valor de produção lançado em \textbf{Configurações > Dados de produção}.\par
\textbf{Como resolver:} Acesse \textbf{Configurações > Dados de produção} e registre o valor produzido daquele mês.
\end{remark}

\section{Visão por Quantidade de Peças e Comparação Temporal}

\subsection*{Objetivo}

Acompanhar o volume físico de itens descartados ou retrabalhados em \textbf{unidades} e escolher qual período passado será usado como referência comparativa no painel.\par

\subsection*{Antes de começar}

\begin{itemize}
\item Estar na tela do \textbf{Dashboard}.
\end{itemize}

\passo[0.96707/0.06966/0.96341/0.01910,0.17805/0.20449/0.13415/0.14045,0.49878/0.20449/0.47561/0.14045,0.81707/0.20449/0.84146/0.14045,0.35366/0.63483/0.35366/0.63483,0.82439/0.63483/0.82439/0.63483]{completo/31-dashboard-quantidade}{Consultar a quantidade de peças}{%
\begin{elementos}
\item \textbf{Quantidade.} Ative a chave para consultar os indicadores em unidades.
\item \textbf{Acumulado.} Confira o total de peças do período.
\item \textbf{Quantidade de referência.} Compare com a quantidade do período de referência.
\item \textbf{Variação da QTY Scrap.} Leia a diferença percentual entre os períodos.
\item \textbf{Gráfico semanal.} Acompanhe a evolução de W01 a W52.
\item \textbf{Produtos/áreas.} Compare a distribuição da quantidade no ranking.
\end{elementos}
}

\passo[0.59070/0.22414/0.79070/0.24138]{completo/32-dashboard-periodo}{Escolher o período de comparação}{
\begin{elementos}
\item \textbf{Período de comparação.} Abra o seletor vs. YTD ant. e escolha YTD anterior (mesmo acumulado do ano anterior), Mesmo mês ano anterior ou Mês anterior.
\end{elementos}
}

\Needspace{160pt}

\subsection*{Campos e controles da tela}

\begingroup
\fontsize{9}{12}\selectfont
\setlength{\LTpre}{4pt}
\setlength{\LTpost}{8pt}
\begin{longtable}{@{}L{3.389cm}L{1.916cm}L{2.652cm}L{6.778cm}@{}}
\toprule
\guiacabtabela \thead{Controle} & \thead{Obrigatório} & \thead{Formato} & \thead{Descrição e funcionamento} \\
\midrule
\endfirsthead
\toprule
\guiacabtabela \thead{Controle} & \thead{Obrigatório} & \thead{Formato} & \thead{Descrição e funcionamento} \\
\midrule
\endhead
\bottomrule
\endlastfoot
\rowcolor{white} \textbf{Chave ``Quantidade''} & Não & Chave de alternância & Altera a exibição dos cartões e gráficos de dólares (\tech{US\$}) para quantidade de peças (\tech{unidades}). \\
\rowcolor{guiaareia!70} \textbf{Período de comparação (``vs. YTD ant.'')} & Sim & Campo de seleção & Define qual período anterior serve de base para calcular as variações percentuais. \\
\end{longtable}
\endgroup

\subsection*{Mensagens e solução de problemas}

\begin{remark}
\textbf{Os gráficos pararam de mostrar valores em dólares (``US\$'')}\par
\textbf{Causa provável:} A chave \textbf{Quantidade} no canto superior direito está ligada.\par
\textbf{Como resolver:} Desligue a chave \textbf{Quantidade} para voltar a ver os valores em dólares (\tech{US\$}).
\end{remark}

\section{Rankings Top 10 e Filtros do Dashboard}

\subsection*{Objetivo}

Visualizar rapidamente os 10 maiores focos de perda da fábrica (por linhas, modelos, componentes e contas ofensoras) e aplicar filtros de período, catálogo ou operação.\par

\subsection*{Antes de começar}

\begin{itemize}
\item Estar na tela do \textbf{Dashboard}.
\end{itemize}

\begin{remark}[info]
A localização e a leitura dos rankings Top 10 já estão descritas no guia rápido. Consulte o passo~\ref{passo:4} (\pagina{passo:4}).
\end{remark}
Os rankings agrupam o impacto por \textbf{linha}, \textbf{modelo}, \textbf{componente}
e \textbf{ofensor}. Exemplos de identificadores exibidos são \tech{A02}, \tech{A05},
\tech{B01}; \tech{Placa principal}, \tech{Moldura superior}, \tech{Tampa traseira};
\tech{Module}, \tech{PCBA}, \tech{Cover Assembly}; e \tech{D-EXPENSE},
\tech{D-DIRECT}, \tech{Intransit Shipment}, \tech{A-STOCK}, respectivamente.


\passo[0.14615/0.38906/0.03846/0.42188,0.46154/0.24063/0.40385/0.17188,0.80000/0.24063/0.84615/0.17188,0.46154/0.41094/0.40385/0.50000,0.80000/0.41094/0.84615/0.50000,0.91923/0.92813/0.92308/0.84375,0.78654/0.92813/0.75000/0.84375]{completo/34-dashboard-filtros}{Filtrar o Dashboard}{%
Abra Filtros no Dashboard ou no gráfico desejado. O detalhe mostra os campos da categoria Operação.\par
\begin{elementos}
\item \textbf{Categorias.} Escolha Período, Catálogo, Operação ou Estrutura.
\item \textbf{Produção.} Restrinja o produto ou o grupo de produção.
\item \textbf{Linha.} Escolha a linha.
\item \textbf{Departamento.} Restrinja o departamento.
\item \textbf{Filtro da linha.} Escolha o agrupamento de linhas.
\item \textbf{Aplicar.} Atualize os gráficos.
\item \textbf{Limpar seleção.} Remova o recorte.
\end{elementos}
}

\Needspace{160pt}

\subsection*{Campos e regras}

\begingroup
\fontsize{9}{12}\selectfont
\setlength{\LTpre}{4pt}
\setlength{\LTpost}{8pt}
\begin{longtable}{@{}L{3.389cm}L{1.916cm}L{2.652cm}L{6.778cm}@{}}
\toprule
\guiacabtabela \thead{Campo} & \thead{Obrigatório} & \thead{Formato} & \thead{Descrição e regra de preenchimento} \\
\midrule
\endfirsthead
\toprule
\guiacabtabela \thead{Campo} & \thead{Obrigatório} & \thead{Formato} & \thead{Descrição e regra de preenchimento} \\
\midrule
\endhead
\bottomrule
\endlastfoot
\rowcolor{white} \textbf{Categorias de Filtro (``Período'', ``Catálogo'', ``Operação'', ``Estrutura'')} & Não & Campos de seleção & Permitem filtrar todo o Dashboard (botão do topo) ou apenas um gráfico individual (botão no canto do próprio gráfico). \\
\end{longtable}
\endgroup

\subsection*{Mensagens e solução de problemas}

\begin{remark}
\textbf{Um gráfico específico apresenta números diferentes dos cartões principais}\par
\textbf{Causa provável:} Há um filtro individual aplicado somente àquele gráfico (o botão \textbf{Filtros} do gráfico fica destacado em rosa).\par
\textbf{Como resolver:} Clique no botão \textbf{Filtros} daquele gráfico e clique em \textbf{Limpar seleção}.
\end{remark}

\fimdecapitulo

\chapter{Relatórios de Análise e Planos de Ação}

\label{cap:manual-6}

\section{Central de Relatórios}

\subsection*{Objetivo}

Listar e organizar todos os relatórios técnicos e planos de ação da fábrica, permitindo acompanhar o status, o andamento das tarefas (\tech{12/12}), a quantidade de análises vinculadas, o custo total do caso e o responsável.\par

\subsection*{Antes de começar}

\begin{itemize}
\item Estar conectado ao sistema e acessar o menu lateral \textbf{Relatórios}.
\end{itemize}

\passo[0.17262/0.20000/0.17262/0.20000,0.96190/0.38182/0.93452/0.30000,0.48214/0.62182/0.48214/0.62182,0.94762/0.06364/0.88690/0.06364]{completo/35-relatorios}{Consultar os relatórios}{%
Abra Relatórios no menu lateral.\par
\begin{elementos}
\item \textbf{Status.} Use Todos, Abertos, Em andamento ou Fechados para localizar os planos.
\item \textbf{Filtros.} Refine a busca pelos relatórios.
\item \textbf{Relatório.} Confira código e título, status, tarefas concluídas/total, análises vinculadas, custo do caso, prazo e responsável. Clique na linha para abrir o relatório.
\item \textbf{Novo relatório.} Inicie o cadastro de um plano.
\end{elementos}
}

\Needspace{160pt}

\subsection*{Campos e regras}

\begingroup
\fontsize{9}{12}\selectfont
\setlength{\LTpre}{4pt}
\setlength{\LTpost}{8pt}
\begin{longtable}{@{}L{3.389cm}L{1.916cm}L{2.652cm}L{6.778cm}@{}}
\toprule
\guiacabtabela \thead{Campo} & \thead{Obrigatório} & \thead{Formato} & \thead{Descrição e regra de preenchimento} \\
\midrule
\endfirsthead
\toprule
\guiacabtabela \thead{Campo} & \thead{Obrigatório} & \thead{Formato} & \thead{Descrição e regra de preenchimento} \\
\midrule
\endhead
\bottomrule
\endlastfoot
\rowcolor{white} \textbf{Abas de Status (``Todos'', ``Abertos'', ``Em andamento'', ``Fechados'')} & Não & Filtro rápido & Filtra a lista de relatórios de acordo com a situação atual de cada plano. \\
\rowcolor{guiaareia!70} \textbf{CUSTO DO CASO} & Somente leitura & Valor em dólares (\tech{US\$}) & Soma automática dos custos de todas as análises de scrap vinculadas ao relatório. \\
\end{longtable}
\endgroup

\subsection*{Mensagens e solução de problemas}

\begin{remark}
\textbf{Um relatório existente não aparece na lista}\par
\textbf{Causa provável:} Alguma aba de status específica (como \tech{Abertos} ou \tech{Fechados}) está selecionada no topo.\par
\textbf{Como resolver:} Clique na aba \textbf{Todos} para exibir todos os relatórios cadastrados.
\end{remark}

\section{Criação de Novo Relatório (Passo a Passo em 4 Etapas)}

\subsection*{Objetivo}

Guiar a criação de um novo relatório técnico em 4 etapas simples: escolha das análises de scrap revisadas, preenchimento dos dados básicos, investigação da causa raiz pelos \textbf{5 Porquês} e registro opcional das primeiras ações.\par

\subsection*{Antes de começar}

\begin{itemize}
\item Existir ao menos \textbf{uma análise de scrap finalizada (``Revisado'')} na Base de Scrap.
\end{itemize}

\passo[0.49839/0.36413/0.09646/0.33696,0.17846/0.44674/0.09646/0.48913,0.92283/0.93478/0.87621/0.88043]{completo/36-novo-relatorio-analises}{Selecionar as análises do relatório}{%
Em Relatórios, use Novo relatório para abrir a etapa 1 de 4.\par
\begin{elementos}
\item \textbf{Busca.} Localize uma análise revisada pelo título, tipo de scrap ou item.
\item \textbf{Selecionar análises.} Marque pelo menos uma ocorrência. Os itens selecionados ficam destacados em rosa.
\item \textbf{Continuar.} Avance para os dados do relatório.
\end{elementos}
}

\passo[0.49839/0.40000/0.09646/0.36957,0.49839/0.56522/0.09646/0.57609,0.27331/0.70326/0.24116/0.77174,0.50000/0.70326/0.54662/0.77174,0.72347/0.70326/0.86817/0.77174,0.92283/0.93804/0.87621/0.88043]{completo/38-novo-relatorio-dados}{Preencher os dados do relatório}{%
\begin{elementos}
\item \textbf{Título.} Informe um título de até 240 caracteres.
\item \textbf{Descrição.} Resuma o problema do relatório.
\item \textbf{Responsável.} Selecione quem conduzirá o plano.
\item \textbf{Prazo.} Informe a data de conclusão.
\item \textbf{Etapa do plano.} Selecione a etapa atual, como Mapeamento de causas.
\item \textbf{Continuar.} Avance para a análise da causa raiz.
\end{elementos}
}

\passo[0.49839/0.56522/0.09646/0.54348,0.49839/0.81522/0.09646/0.81522,0.92283/0.95109/0.87621/0.89674]{completo/39-novo-relatorio-5-porques}{Preencher os 5 Porquês}{%
\begin{elementos}
\item \textbf{Porquês 1 a 3.} Preencha obrigatoriamente os três primeiros níveis, perguntando por que o problema ocorreu até identificar a causa raiz.
\item \textbf{Porquês 4 e 5.} Use os campos opcionais para aprofundar a investigação, quando necessário.
\item \textbf{Continuar.} Avance para as ações iniciais.
\end{elementos}
}

\passo[0.23955/0.49239/0.12862/0.56522,0.90997/0.94022/0.86817/0.89130]{completo/40-novo-relatorio-acoes}{Adicionar ações iniciais}{
\begin{elementos}
\item \textbf{Adicionar ação.} Na etapa 4 de 4, cadastre as primeiras tarefas, se desejar. Você também poderá incluí-las depois, no relatório.
\item \textbf{Criar relatório.} Conclua o cadastro com ou sem ações iniciais.
\end{elementos}
}

\Needspace{160pt}

\subsection*{Campos e regras}

\begingroup
\fontsize{9}{12}\selectfont
\setlength{\LTpre}{4pt}
\setlength{\LTpost}{8pt}
\begin{longtable}{@{}L{3.389cm}L{1.916cm}L{2.652cm}L{6.778cm}@{}}
\toprule
\guiacabtabela \thead{Campo} & \thead{Obrigatório} & \thead{Formato} & \thead{Descrição e regra de preenchimento} \\
\midrule
\endfirsthead
\toprule
\guiacabtabela \thead{Campo} & \thead{Obrigatório} & \thead{Formato} & \thead{Descrição e regra de preenchimento} \\
\midrule
\endhead
\bottomrule
\endlastfoot
\rowcolor{white} \textbf{Análises vinculadas (Etapa 1)} & Sim (mínimo 1) & Caixas de marcação & Vincula ao relatório uma ou mais ocorrências já revisadas na Base de Scrap. \\
\rowcolor{guiaareia!70} \textbf{Título * (Etapa 2)} & Sim & Texto (até 240 caracteres) & Nome principal do relatório e do plano de ação. \\
\rowcolor{white} \textbf{Descrição (Etapa 2)} & Não & Texto livre & Resumo do contexto e objetivo do relatório. \\
\rowcolor{guiaareia!70} \textbf{Responsável / Prazo (Etapa 2)} & Não & Seletor e data (\tech{AAAA/MM/DD}) & Definem quem responde pelo plano e a data-limite. A etapa do plano possui uma opção inicial que pode ser ajustada no seletor. \\
\rowcolor{white} \textbf{Porquê 1, Porquê 2 e Porquê 3 (Etapa 3)} & Sim & Texto livre & Os 3 primeiros porquês são obrigatórios para registrar a cadeia de causa raiz. \\
\rowcolor{guiaareia!70} \textbf{Porquê 4 e Porquê 5 (Etapa 3)} & Não & Texto livre & Campos complementares opcionais da análise de 5 Porquês. \\
\rowcolor{white} \textbf{Ações iniciais (Etapa 4)} & Não & Lista de tarefas & Permite adiantar o cadastro de ações corretivas antes de abrir o relatório. \\
\end{longtable}
\endgroup

\subsection*{Mensagens e solução de problemas}

\begin{remark}
\textbf{Botão "Continuar" bloqueado na Etapa 1}\par
\textbf{Causa provável:} Nenhuma análise de scrap foi marcada na lista.\par
\textbf{Como resolver:} Marque pelo menos uma análise revisada para liberar o avanço.
\end{remark}

\begin{remark}
\textbf{Botão "Continuar" bloqueado na Etapa 3}\par
\textbf{Causa provável:} Algum dos campos obrigatórios (\textbf{Porquê 1}, \textbf{Porquê 2} ou \textbf{Porquê 3}) está em branco.\par
\textbf{Como resolver:} Preencha no mínimo os 3 primeiros motivos da investigação de causa.
\end{remark}

\section{Abas Visão geral e Investigação do Relatório}

\subsection*{Objetivo}

Acompanhar na aba \textbf{Visão geral} os indicadores financeiros e a lista de \textbf{Prontidão do relatório} (sete itens de verificação do preenchimento) e detalhar na aba \textbf{Investigação} a causa raiz, os 5 porquês e as análises vinculadas.\par

\subsection*{Antes de começar}

\begin{itemize}
\item Estar com um relatório aberto na tela de \textbf{Relatórios}.
\end{itemize}

\passo[0.09036/0.19178/0.08434/0.10959,0.32771/0.24795/0.39157/0.31507,0.54578/0.24795/0.60843/0.31507,0.76506/0.19178/0.88554/0.10959,0.24699/0.61644/0.24699/0.61644]{completo/41-relatorio-visao-geral}{Consultar a visão geral}{%
\begin{elementos}
\item \textbf{Custo do caso.} Confira o custo consolidado das análises vinculadas.
\item \textbf{Economia estimada.} Use Informar economia para registrar o valor esperado.
\item \textbf{Redução observada.} Use Ver efeito para consultar a medição na base.
\item \textbf{Evidências pendentes.} Confira o que falta registrar.
\item \textbf{Prontidão do relatório.} Clique no item pendente para ir ao seu preenchimento: causa raiz, 4M, tarefa, medição antes/depois, mês do efeito ou resultados qualitativos.
\end{elementos}
}

\passo[0.11257/0.08304/0.08383/0.04678,0.94371/0.19298/0.89222/0.15789,0.93293/0.34971/0.88024/0.40936,0.35329/0.74854/0.35329/0.74854,0.90898/0.75439/0.85629/0.83041]{completo/42-relatorio-investigacao}{Consultar a investigação}{%
\begin{elementos}
\item \textbf{Investigação.} Abra a aba do relatório.
\item \textbf{Editar causa raiz.} Registre a definição final da causa.
\item \textbf{Editar plano.} Revise a análise dos 5 Porquês.
\item \textbf{Análises vinculadas.} Confira as ocorrências que compõem o relatório.
\item \textbf{Caso BP.} Destaque uma ocorrência como boa prática para a apresentação exportada.
\end{elementos}
}

\Needspace{160pt}

\subsection*{Campos e regras}

\begingroup
\fontsize{9}{12}\selectfont
\setlength{\LTpre}{4pt}
\setlength{\LTpost}{8pt}
\begin{longtable}{@{}L{3.389cm}L{1.916cm}L{2.652cm}L{6.778cm}@{}}
\toprule
\guiacabtabela \thead{Campo} & \thead{Obrigatório} & \thead{Formato} & \thead{Descrição e regra de preenchimento} \\
\midrule
\endfirsthead
\toprule
\guiacabtabela \thead{Campo} & \thead{Obrigatório} & \thead{Formato} & \thead{Descrição e regra de preenchimento} \\
\midrule
\endhead
\bottomrule
\endlastfoot
\rowcolor{white} \textbf{Prontidão do relatório (``1 de 7'' a ``7 de 7'')} & Somente leitura & Lista de verificação & Indica quais informações ainda faltam para completar o relatório. A lista é informativa e não impede salvar ou exportar o documento. \\
\rowcolor{guiaareia!70} \textbf{Causa raiz} & Recomendado & Texto livre & Descrição definitiva da causa fundamental do problema. \\
\rowcolor{white} \textbf{Botão ``\ensuremath{\heartsuit} Caso BP''} & Não & Botão de destaque & Marca a análise selecionada como estudo de caso em destaque para exportação em PDF e PowerPoint. \\
\end{longtable}
\endgroup

\subsection*{Mensagens e solução de problemas}

\begin{remark}
\textbf{Itens pendentes na lista de Prontidão do relatório}\par
\textbf{Causa provável:} Informações de causa raiz, tarefas ou medição de resultado ainda não foram salvas nas abas correspondentes.\par
\textbf{Como resolver:} Clique sobre cada item desmarcado na lista \textbf{Prontidão do relatório} para abrir e preencher a seção faltante.
\end{remark}

\Needspace{190pt}
\section{Gestão de Tarefas (Abas Lista, Quadro e Cronograma)}

\subsection*{Objetivo}

Cadastrar, atribuir responsáveis, acompanhar prazos e validar a conclusão das ações corretivas, preventivas e de contenção do relatório em três formatos de visualização: \textbf{Lista}, \textbf{Quadro} de colunas e \textbf{Cronograma} no tempo.\par

\subsection*{Antes de começar}

\begin{itemize}
\item Estar dentro de um relatório aberto.
\end{itemize}

\passo[0.16407/0.07534/0.12575/0.02055,0.52096/0.22329/0.52096/0.22329,0.35329/0.74521/0.36527/0.83562,0.95329/0.74521/0.92814/0.83562]{completo/43-tarefas-lista}{Acompanhar tarefas em lista}{%
\begin{elementos}
\item \textbf{Lista.} Consulte as tarefas em formato de tabela.
\item \textbf{Filtros.} Localize tarefas por busca, estado, tipo, responsável ou agrupamento.
\item \textbf{Título da nova tarefa.} Digite a descrição curta da ação no campo inferior.
\item \textbf{Adicionar.} Inclua a tarefa. A tabela apresenta departamento, 4M, responsável, estado, prioridade, prazo e os demais dados de acompanhamento.
\end{elementos}
}

\passo[0.21557/0.07853/0.17365/0.02454,0.13413/0.25399/0.13413/0.25399,0.39401/0.25399/0.39401/0.25399,0.64192/0.25399/0.64192/0.25399,0.89222/0.25399/0.89222/0.25399]{completo/44-tarefas-quadro}{Acompanhar tarefas no quadro}{%
\begin{elementos}
\item \textbf{Quadro.} Abra a visão de cartões.
\item \textbf{Planejada.} Identifique as ações ainda não iniciadas.
\item \textbf{Em andamento.} Acompanhe as ações em execução.
\item \textbf{Em validação.} Localize as tarefas aguardando aprovação.
\item \textbf{Concluída.} Confira as ações finalizadas. Para cadastrar uma tarefa diretamente na primeira coluna, consulte o passo~\ref{passo:37}.
\end{elementos}
}

\passo[0.48347/0.26667/0.44215/0.11111,0.11157/0.61333/0.09504/0.86667,0.23967/0.61333/0.27686/0.86667,0.32645/0.61333/0.45041/0.86667,0.85950/0.61333/0.81405/0.86667]{completo/45-tarefa-criacao-rapida}{Criar uma tarefa rapidamente}{%
No Quadro, use Criar tarefa na coluna Planejada para abrir o formulário rápido.\par
\begin{elementos}
\item \textbf{Título.} Em Criar tarefa, descreva o que precisa ser feito.
\item \textbf{Prioridade.} Escolha a prioridade pelo ícone correspondente.
\item \textbf{Prazo.} Defina a data no ícone de calendário.
\item \textbf{Responsável.} Selecione o responsável pela tarefa.
\item \textbf{Confirmar.} Pressione Enter ou use o ícone de confirmação para criar o cartão na coluna Planejada.
\end{elementos}
}

\passo[0.44000/0.84828/0.76000/0.96552]{completo/46-tarefa-validar-conclusao}{Validar a conclusão de tarefa}{%
\begin{elementos}
\item \textbf{Validar conclusão.} Com perfil Gestor, clique em Validar conclusão na tarefa que estiver na coluna Em validação. O cartão será movido para Concluída.
\end{elementos}
}

\passo[0.49620/0.24571/0.05894/0.20000,0.49620/0.42857/0.05894/0.44000,0.38783/0.59314/0.19202/0.65714,0.86502/0.59314/0.91445/0.65714,0.49620/0.78171/0.32510/0.82857,0.05513/0.84800/0.04943/0.91429,0.52662/0.84800/0.59125/0.91429]{completo/47-nova-tarefa-dados}{Cadastrar os dados da tarefa}{%
\begin{elementos}
\item \textbf{Título.} Use Nova tarefa no relatório e informe o título da ação.
\item \textbf{Descrição.} Explique o trabalho a executar.
\item \textbf{Tipo de ação.} Selecione Não classificada, Contenção, Corretiva ou Preventiva.
\item \textbf{Prioridade.} Defina o nível de prioridade.
\item \textbf{Reaproveitar dados.} Aproveite a classificação de uma análise vinculada, quando aplicável.
\item \textbf{4M.} Role um pouco para preencher a categoria completa.
\item \textbf{Sintoma.} Preencha o sintoma ou mantenha o valor reaproveitado da análise.
\end{elementos}
\par\smallskip {\fontsize{8.2}{11}\selectfont\color{guiacinza}Os rótulos de 4M e Sintoma aparecem no limite inferior da captura. Role a janela para mostrar esses campos por inteiro e prossiga para responsável e prazo, no passo~\ref{passo:40}.\par}
}

\passo[0.27875/0.33708/0.06433/0.29213,0.76608/0.33708/0.90253/0.29213,0.27875/0.45618/0.06433/0.42135,0.76608/0.45618/0.90253/0.42135,0.51267/0.55618/0.51267/0.59551,0.14230/0.62921/0.08382/0.68539,0.51267/0.72135/0.59064/0.73034,0.51267/0.83146/0.59064/0.85393,0.92203/0.94382/0.88304/0.88764]{completo/48-nova-tarefa-responsavel-prazo}{Definir responsável, prazo e economia}{%
\begin{elementos}
\item \textbf{Responsável principal.} Selecione quem executará a tarefa.
\item \textbf{Área responsável.} Informe a área, como Qualidade.
\item \textbf{Início.} Defina a data de início.
\item \textbf{Prazo.} Informe a data de entrega.
\item \textbf{Participantes.} Inclua as pessoas envolvidas.
\item \textbf{Bloqueio.} Indique Sem bloqueio ou Bloqueada.
\item \textbf{Implantação por linha.} Expanda para indicar as linhas de produção atendidas.
\item \textbf{Economia.} Expanda para informar o ganho financeiro esperado.
\item \textbf{Salvar.} Grave a tarefa.
\end{elementos}
}

\passo[0.45122/0.47273/0.45122/0.47273,0.94390/0.05455/0.90244/0.10303]{completo/50-tarefas-cronograma}{Acompanhar o cronograma}{
\begin{elementos}
\item \textbf{Cronograma das tarefas.} Abra a aba Cronograma e acompanhe as barras de duração das ações ao longo dos meses.
\item \textbf{Ocultar concluídas.} Marque a opção para exibir apenas as tarefas ainda em aberto.
\end{elementos}
}

\Needspace{160pt}

\subsection*{Campos e regras}

\begingroup
\fontsize{9}{12}\selectfont
\setlength{\LTpre}{4pt}
\setlength{\LTpost}{8pt}
\begin{longtable}{@{}L{3.389cm}L{1.916cm}L{2.652cm}L{6.778cm}@{}}
\toprule
\guiacabtabela \thead{Campo} & \thead{Obrigatório} & \thead{Formato} & \thead{Descrição e regra de preenchimento} \\
\midrule
\endfirsthead
\toprule
\guiacabtabela \thead{Campo} & \thead{Obrigatório} & \thead{Formato} & \thead{Descrição e regra de preenchimento} \\
\midrule
\endhead
\bottomrule
\endlastfoot
\rowcolor{white} \textbf{Título * (da Tarefa)} & Sim & Texto curto & Nome da ação prática que deverá ser executada. \\
\rowcolor{guiaareia!70} \textbf{Tipo de ação} & Sim & Botões de opção (\tech{Não classificada}, \tech{Contenção}, \tech{Corretiva}, \tech{Preventiva}) & Define a finalidade técnica da tarefa no plano de ação. \\
\rowcolor{white} \textbf{Prioridade} & Sim & Campo de seleção & Define o nível de urgência da tarefa (\tech{Baixa}, \tech{Média}, \tech{Alta}, \tech{Crítica}). \\
\rowcolor{guiaareia!70} \textbf{Responsável principal / Área responsável} & Não & Seletor de colaborador e Texto & Define quem executará a tarefa (enviando alerta automático para o responsável) e o setor correspondente. \\
\rowcolor{white} \textbf{Início / Prazo} & Não & Datas (\tech{AAAA/MM/DD}) & Datas previstas para a tarefa, exibidas no \textbf{Cronograma}. Quando ambas forem informadas, o início não pode ser posterior ao prazo. \\
\rowcolor{guiaareia!70} \textbf{Situação de bloqueio (``Sem bloqueio'' / ``Bloqueada'')} & Não & Botões de opção & Indica se há algum impedimento que impede o andamento da ação. \\
\rowcolor{white} \textbf{Implantação por linha / Economia} & Não & Seções expansíveis & Registram em quais linhas a melhoria foi aplicada e a economia anual estimada. \\
\end{longtable}
\endgroup

\subsection*{Mensagens e solução de problemas}

\begin{remark}
\textbf{Data em vermelho com símbolo de alerta (``\faExclamationTriangle{}'') no cartão da tarefa}\par
\textbf{Causa provável:} O prazo final da tarefa já passou e ela ainda não foi concluída.\par
\textbf{Como resolver:} Atualize o andamento junto ao responsável ou edite a tarefa para ajustar o prazo caso tenha sido reprogramada.
\end{remark}

\section{Resultado, prévia e exportação de relatórios}

\subsection*{Objetivo}

Consultar na aba \textbf{Resultado} a variação de scrap após as ações (comparando os 3 meses anteriores com os 3 meses posteriores ao mês de efeito), revisar o resumo na aba \textbf{Prévia} e baixar o documento pronto em \textbf{PDF A4} ou \textbf{Apresentação PowerPoint 16:9}.\par

\subsection*{Antes de começar}

\begin{itemize}
\item Análises e plano de ação registrados no relatório.
\end{itemize}

\passo[0.30864/0.13533/0.20988/0.11377,0.05679/0.21078/0.13580/0.22754,0.25926/0.32814/0.17284/0.38323,0.41975/0.63473/0.41975/0.63473,0.85926/0.63473/0.85926/0.63473,0.23457/0.95449/0.23457/0.95449]{completo/51-relatorio-resultado}{Consultar os resultados}{%
Abra Resultado no relatório. A captura mostra a consulta dos indicadores pós-implantação.\par
\begin{elementos}
\item \textbf{Medição antes e depois.} Leia o resumo dos resultados comparados.
\item \textbf{Mês do efeito.} Confira o mês usado como referência da implantação.
\item \textbf{Resultados qualitativos.} Leia as melhorias descritas para o caso.
\item \textbf{Gráfico antes/depois.} Compare os três meses anteriores, em cinza, e os três seguintes, em rosa.
\item \textbf{Resumo do efeito.} Confira os custos antes/depois e as variações financeira e de quantidade.
\item \textbf{Por análise.} Consulte a comparação por ocorrência vinculada.
\end{elementos}
}

\passo[0.40244/0.36571/0.40244/0.36571,0.46341/0.84800/0.46341/0.84800]{completo/52-relatorio-fotos}{Conferir as fotos da melhoria}{
Na aba Resultado, role a página até as evidências visuais da melhoria.\par
\begin{elementos}
\item \textbf{Fotos antes e depois.} Compare a situação inicial com o resultado das ações. O painel apresenta até seis pares de imagens selecionados entre os anexos das ações do relatório.
\item \textbf{Fotos do processo.} Confira os registros de Mapeamento de causas, Geração de ideias, Execução e Gestão de resultados.
\end{elementos}
}

\passo[0.42424/0.50568/0.42424/0.50568,0.95273/0.17045/0.94545/0.28977]{completo/53-relatorio-previa}{Revisar a prévia do relatório}{
\begin{elementos}
\item \textbf{Prévia do relatório.} Abra a aba Prévia e revise o conteúdo consolidado antes de exportar.
\item \textbf{Menu de ações.} Use o botão de três pontos para acessar Ver detalhes, Editar plano ou Excluir relatório, conforme a ação desejada.
\end{elementos}
}

\passo[0.62500/0.17354/0.67029/0.13228,0.18841/0.30899/0.09058/0.28042,0.14130/0.43069/0.07246/0.46032,0.43478/0.43386/0.50725/0.48148,0.39855/0.64339/0.39855/0.64339,0.84964/0.55556/0.84964/0.55556,0.89674/0.94709/0.92391/0.89418]{completo/54-exportar-pdf}{Exportar o relatório em PDF}{%
Na Prévia do relatório, use Baixar documento para abrir Montar documento.\par
\begin{elementos}
\item \textbf{Tradução por IA.} Ative Traduzir o texto para inglês quando precisar traduzir automaticamente o texto digitado.
\item \textbf{PDF A4.} Escolha o formato executivo para ler e imprimir.
\item \textbf{Idioma.} Selecione Português, English ou 한국어 para os títulos e rótulos do documento.
\item \textbf{Atalhos.} Use Completo, Só visão executiva ou Só casos como ponto de partida.
\item \textbf{Seções.} Marque quais blocos entrarão no arquivo: ficha anual do programa (Task Charter), resumo, efeito na base, casos e gráficos.
\item \textbf{Estrutura.} Confira a ordem e a prévia das páginas.
\item \textbf{Baixar PDF.} Gere o arquivo final.
\end{elementos}
}

\passo[0.50725/0.30899/0.59783/0.28042,0.14130/0.43069/0.07246/0.46032,0.85145/0.55556/0.85145/0.55556,0.89674/0.94709/0.92391/0.89418]{completo/55-exportar-powerpoint}{Exportar a apresentação}{%
\begin{elementos}
\item \textbf{Apresentação 16:9.} Selecione os slides editáveis no PowerPoint.
\item \textbf{Idioma.} Escolha Português, English ou 한국어.
\item \textbf{Estrutura.} Confira a lista e a ordem dos slides.
\item \textbf{Baixar apresentação.} Gere o arquivo para revisar no PowerPoint.
\end{elementos}
}

\passo[0.12717/0.29834/0.12717/0.29834,0.13873/0.51934/0.13873/0.51934,0.47399/0.51934/0.47399/0.51934,0.79769/0.30939/0.79769/0.30939,0.79769/0.71823/0.79769/0.71823]{completo/56-slide-exportado}{Conferir o slide exportado}{%
\begin{elementos}
\item \textbf{Problema.} Confira o título do caso, os dados da ocorrência e o custo.
\item \textbf{Causa raiz.} Revise a causa e a classificação apresentadas.
\item \textbf{Antes e depois.} Confira as fotos e suas descrições.
\item \textbf{Ações.} Revise as ações executadas, responsáveis, datas e status.
\item \textbf{Resultados.} Confira a redução de custo, a evolução dos indicadores e a economia prevista antes de compartilhar os slides.
\end{elementos}
}

\Needspace{160pt}

\subsection*{Campos e regras}

\begingroup
\fontsize{9}{12}\selectfont
\setlength{\LTpre}{4pt}
\setlength{\LTpost}{8pt}
\begin{longtable}{@{}L{3.389cm}L{1.916cm}L{2.652cm}L{6.778cm}@{}}
\toprule
\guiacabtabela \thead{Campo} & \thead{Obrigatório} & \thead{Formato} & \thead{Descrição e regra de preenchimento} \\
\midrule
\endfirsthead
\toprule
\guiacabtabela \thead{Campo} & \thead{Obrigatório} & \thead{Formato} & \thead{Descrição e regra de preenchimento} \\
\midrule
\endhead
\bottomrule
\endlastfoot
\rowcolor{white} \textbf{Mês do efeito (aba ``Resultado'')} & Para medir o efeito & Mês e Ano (\tech{mmm. AAAA}) & Mês de referência da melhoria, necessário para comparar os três meses anteriores com os três posteriores. \\
\rowcolor{guiaareia!70} \textbf{Traduzir o texto para inglês (``IA'')} & Não & Chave liga/desliga & Traduz automaticamente para o inglês os textos abertos do relatório antes de gerar o arquivo. \\
\rowcolor{white} \textbf{Formato (``PDF A4'' / ``Apresentação 16:9'')} & Sim & Seleção visual & Escolhe entre gerar um documento vertical em \tech{.pdf} ou slides widescreen editáveis em \tech{.pptx}. \\
\rowcolor{guiaareia!70} \textbf{Idioma do documento} & Sim & Campo de seleção (\tech{Português}, \tech{English}, {\color{guiavermelho}한국어}) & Define o idioma dos títulos e rótulos do arquivo gerado. \\
\rowcolor{white} \textbf{Atalhos e Chaves de Seção} & Não & Botões e Chaves liga/desliga & Permitem escolher exatamente quais blocos e gráficos entram no documento exportado. \\
\end{longtable}
\endgroup

\subsection*{Mensagens e solução de problemas}

\begin{remark}
\textbf{Taxa (scrap ÷ produção) exibe "sem produção cadastrada" na aba Resultado}\par
\textbf{Causa provável:} Os meses usados na comparação antes/depois ainda não possuem dados de produção salvos em \textbf{Configurações > Dados de produção}.\par
\textbf{Como resolver:} Cadastre o volume produzido desses meses em \textbf{Configurações > Dados de produção} para que o sistema calcule a taxa relativa.
\end{remark}

\fimdecapitulo

\part{Administração e Recursos Complementares}
\chapter{Governança de Alertas e Integração com Microsoft Teams}

\label{cap:manual-7}

\section{Histórico de Notificações e Regras de Alerta}

\subsection*{Objetivo}

Consultar os avisos recebidos pelo operador (como novas tarefas atribuídas, prazos vencidos ou limites de custo ultrapassados) e configurar regras automáticas para disparo de alertas no sistema e no Microsoft Teams.\par

\subsection*{Antes de começar}

\begin{itemize}
\item Estar conectado ao sistema e acessar o menu lateral \textbf{Alertas}.
\end{itemize}

\passo[0.43976/0.46328/0.43976/0.46328,0.85663/0.30169/0.81325/0.24859,0.91928/0.30169/0.90964/0.22599,0.94819/0.30169/0.97711/0.36158,0.92651/0.03955/0.88554/0.10169]{completo/57-alertas}{Consultar o histórico de alertas}{%
\begin{elementos}
\item \textbf{Histórico.} Abra Alertas e confira título, classificação, data, status e responsável.
\item \textbf{Ver tarefa.} Abra a tarefa associada ao aviso.
\item \textbf{Editar.} Use o lápis na mesma linha para editar.
\item \textbf{Marcar como lido.} Use o ícone de confirmação para registrar a leitura.
\item \textbf{Configurar alertas.} Abra o cadastro de regras, detalhado no passo~\ref{passo:49}.
\end{elementos}
}

\passo[0.38889/0.53488/0.38889/0.53488,0.74722/0.13721/0.63889/0.06977,0.85000/0.40930/0.88889/0.27907]{completo/58-configurar-alertas}{Configurar os alertas}{%
\begin{elementos}
\item \textbf{Regras cadastradas.} Em Configurar alertas, localize a regra desejada.
\item \textbf{Nova regra.} Inicie uma regra de aviso.
\item \textbf{Editar.} Altere uma regra existente pelo botão da linha.
\end{elementos}
}

\passo[0.50000/0.21571/0.06522/0.19372,0.26087/0.31309/0.06522/0.29319,0.71739/0.31309/0.89130/0.29319,0.06522/0.40314/0.10870/0.45026,0.47826/0.57173/0.08696/0.54450,0.25217/0.69738/0.06522/0.67016,0.71957/0.69738/0.89130/0.67016,0.25217/0.84188/0.08696/0.81675,0.71957/0.84188/0.89130/0.81675,0.19565/0.89529/0.60870/0.90890]{completo/59-regra-alerta-condicoes}{Definir quando avisar}{%
\begin{elementos}
\item \textbf{Nome da regra.} Identifique claramente o aviso.
\item \textbf{Tipo.} Escolha a condição, como Custo acima do limite.
\item \textbf{Severidade.} Defina a classificação do alerta.
\item \textbf{Habilitada.} Mantenha ligada para avaliar a regra.
\item \textbf{Avaliar custo por.} Escolha Linha ou Ocorrência.
\item \textbf{Limiar.} Informe o custo que dispara o aviso, por exemplo 4.000.
\item \textbf{Moeda.} Selecione a moeda do limite.
\item \textbf{Início do período.} Informe a data inicial da avaliação.
\item \textbf{Fim do período.} Informe a data final.
\item \textbf{Intervalo mínimo.} Role um pouco para preencher o intervalo, em minutos, entre avisos.
\end{elementos}
\par\smallskip {\fontsize{8.2}{11}\selectfont\color{guiacinza}O último marcador identifica o rótulo Intervalo mínimo; a caixa de valor continua abaixo do rodapé mostrado. Complete os destinatários e os canais no passo~\ref{passo:51}.\par}
}

\passo[0.51000/0.20928/0.51000/0.20928,0.51000/0.56701/0.11000/0.48454,0.11000/0.62474/0.21000/0.65464,0.11400/0.73608/0.11000/0.79381,0.55000/0.73608/0.57000/0.79381,0.88800/0.94021/0.87000/0.88660]{completo/60-regra-alerta-destinatarios}{Definir destinatários e canais}{%
\begin{elementos}
\item \textbf{Filtros de escopo.} Se necessário, restrinja Organização, Linha, Part number ou Componente/tipo de item.
\item \textbf{Pessoas.} Selecione ao menos uma pessoa e/ou um perfil destinatário.
\item \textbf{Perfis.} Selecione os perfis que devem receber o aviso, quando aplicável.
\item \textbf{No sistema.} Mantenha o canal interno marcado.
\item \textbf{Microsoft Teams.} Marque também este canal se o destino de Teams estiver configurado.
\item \textbf{Salvar regra.} Grave as condições e os destinatários.
\end{elementos}
}

\Needspace{160pt}

\subsection*{Campos e regras}

\begingroup
\fontsize{9}{12}\selectfont
\setlength{\LTpre}{4pt}
\setlength{\LTpost}{8pt}
\begin{longtable}{@{}L{3.389cm}L{1.916cm}L{2.652cm}L{6.778cm}@{}}
\toprule
\guiacabtabela \thead{Campo} & \thead{Obrigatório} & \thead{Formato} & \thead{Descrição e regra de preenchimento} \\
\midrule
\endfirsthead
\toprule
\guiacabtabela \thead{Campo} & \thead{Obrigatório} & \thead{Formato} & \thead{Descrição e regra de preenchimento} \\
\midrule
\endhead
\bottomrule
\endlastfoot
\rowcolor{white} \textbf{Nome da regra} & Sim & Texto curto & Nome que identifica o alerta na lista de configuração. \\
\rowcolor{guiaareia!70} \textbf{Tipo / Severidade} & Sim & Campos de seleção & Definem o motivo do disparo (\tech{Custo acima do limite}, \tech{Ocorrência relevante}, etc.) e a gravidade (\tech{Informativo}, \tech{Crítico}). \\
\rowcolor{white} \textbf{Habilitada} & Sim & Chave liga/desliga & Permite pausar temporariamente uma regra sem precisar excluí-la. \\
\rowcolor{guiaareia!70} \textbf{Avaliar custo por} & Sim & Campo de seleção (\tech{Linha}, \tech{Ocorrência}) & Define se o sistema soma o custo acumulado da linha ou avalia cada ocorrência individualmente. \\
\rowcolor{white} \textbf{Limiar / Moeda} & Sim & Número (\tech{>= 0}) e Moeda (\tech{USD}) & Valor financeiro a partir do qual o alerta é disparado. \\
\rowcolor{guiaareia!70} \textbf{Início do período / Fim do período} & Sim & Datas (\tech{AAAA/MM/DD}) & Período considerado na avaliação do custo; a data final deve ser igual ou posterior à data inicial. \\
\rowcolor{white} \textbf{QUEM RECEBE (``Pessoas'' / ``Perfis'')} & Sim (mínimo 1) & Seleção de usuários/perfis & Indica quem receberá a notificação quando a regra for acionada. \\
\rowcolor{guiaareia!70} \textbf{CANAIS DE NOTIFICAÇÃO} & Sim & Caixas de marcação & O envio \textbf{No sistema} é padrão; marque \textbf{Microsoft Teams} para enviar também ao Teams. \\
\end{longtable}
\endgroup

\subsection*{Mensagens e solução de problemas}

\begin{remark}
\textbf{Não é possível salvar a regra de alerta}\par
\textbf{Causa provável:} Falta preencher o nome da regra, escolher pelo menos um destinatário em \textbf{QUEM RECEBE}, ou a data final é anterior à data inicial.\par
\textbf{Como resolver:} Revise esses campos e clique novamente em \textbf{Salvar regra}.
\end{remark}

\section{Integração com Microsoft Teams}

\subsection*{Objetivo}

Cadastrar os endereços de conexão (\emph{Workflows}) do Microsoft Teams para que o sistema Hanaro possa enviar cartões automáticos de alerta para canais, grupos ou pessoas no Teams. Operado pelo perfil \textbf{Administrador}.\par

\subsection*{Antes de começar}

\begin{itemize}
\item Criar previamente um \emph{Workflow} no Microsoft Teams (ou Power Automate) com o gatilho \emph{"When a Teams webhook request is received"} liberado para \emph{Anyone} e copiar o endereço (URL) gerado.
\end{itemize}

\passo[0.42935/0.26744/0.42935/0.26744,0.95000/0.44767/0.90761/0.43023]{completo/61-integracao-teams}{Abrir a integração com Teams}{
Na visão administrativa, abra Integrações pelo ícone de globo no menu lateral.\par
\begin{elementos}
\item \textbf{Orientação da integração.} Leia os requisitos para conectar o Hanaro ao Microsoft Teams.
\item \textbf{Novo destino.} Abra o formulário de cadastro do canal, grupo ou pessoa.
\end{elementos}
}

\passo[0.26011/0.34426/0.11475/0.30601,0.71257/0.34426/0.92350/0.30601,0.26230/0.46776/0.11475/0.43716,0.49727/0.59016/0.11475/0.57377,0.04699/0.67322/0.06011/0.72678,0.51038/0.67322/0.54098/0.72678,0.93005/0.76503/0.91257/0.83060,0.46448/0.06667/0.46448/0.06667]{completo/62-teams-novo-destino-ingles}{Cadastrar um destino no Teams}{%
\begin{elementos}
\item \textbf{Nome do destino.} Informe um nome claro para identificar a pessoa, o grupo ou o canal.
\item \textbf{Tipo.} Escolha Pessoa, Grupo ou Canal.
\item \textbf{Idioma do cartão.} Escolha o idioma do aviso enviado ao Teams, como Português ou Inglês.
\item \textbf{URL do webhook.} Cole o endereço do Workflow do Microsoft Teams.
\item \textbf{Ativo.} Mantenha o destino habilitado para receber avisos.
\item \textbf{Destino operacional padrão.} Marque se o destino receberá por padrão os alertas operacionais gerais.
\item \textbf{Salvar.} Grave o destino.
\item \textbf{Confirmação.} Após salvar, confira a mensagem ``Workflow aceitou a solicitação; confira o cartão no Teams.'' Abra o Teams para confirmar o recebimento.
\end{elementos}
}

\Needspace{160pt}

\subsection*{Campos e regras}

\begingroup
\fontsize{9}{12}\selectfont
\setlength{\LTpre}{4pt}
\setlength{\LTpost}{8pt}
\begin{longtable}{@{}L{3.389cm}L{1.916cm}L{2.652cm}L{6.778cm}@{}}
\toprule
\guiacabtabela \thead{Campo} & \thead{Obrigatório} & \thead{Formato} & \thead{Descrição e regra de preenchimento} \\
\midrule
\endfirsthead
\toprule
\guiacabtabela \thead{Campo} & \thead{Obrigatório} & \thead{Formato} & \thead{Descrição e regra de preenchimento} \\
\midrule
\endhead
\bottomrule
\endlastfoot
\rowcolor{white} \textbf{Nome do destino} & Sim & Texto curto & Nome de identificação do canal, grupo ou colaborador no sistema. \\
\rowcolor{guiaareia!70} \textbf{Tipo} & Sim & Campo de seleção (\tech{Pessoa}, \tech{Grupo}, \tech{Canal}) & Define o tipo de destinatário no Microsoft Teams. \\
\rowcolor{white} \textbf{Idioma do cartão} & Sim & Campo de seleção (\tech{Português}, \tech{Inglês}) & Define o idioma do cartão de aviso enviado ao Teams. \\
\rowcolor{guiaareia!70} \textbf{URL do webhook} & Sim & Endereço web (\tech{https://...}) & Link gerado pelo \emph{Workflow} configurado no Microsoft Teams. \\
\rowcolor{white} \textbf{Ativo / Destino operacional padrão} & Não & Caixas de marcação & Habilitam o envio para este endereço e definem se ele é o canal padrão da operação. \\
\end{longtable}
\endgroup

\subsection*{Mensagens e solução de problemas}

\begin{remark}[info]
\textbf{Workflow aceitou a solicitação; confira o cartão no Teams.}\par
\textbf{O que significa:} O Workflow aceitou a solicitação de envio; é necessário conferir se o cartão chegou ao destino.\par
\textbf{Próximo passo:} Abra o canal ou chat no Microsoft Teams para confirmar o recebimento do cartão de teste.
\end{remark}

\fimdecapitulo

\chapter{Suporte, Gestão de Usuários e Uso da Plataforma}

\label{cap:manual-8}

\section{Abertura de Chamados e Caixa de Entrada de Suporte}

\subsection*{Objetivo}

Permitir que qualquer colaborador envie dúvidas, sugestões ou relatos de problemas diretamente pelo menu \textbf{Ajuda e suporte}, e oferecer aos Administradores uma \textbf{Caixa de entrada} para consultar as solicitações enviadas.\par

\subsection*{Antes de começar}

\begin{itemize}
\item Estar conectado ao sistema.
\end{itemize}

\passo[0.50142/0.35814/0.13314/0.31783,0.50142/0.62791/0.13314/0.59690,0.80737/0.90233/0.78470/0.81395]{completo/65-fale-com-suporte}{Enviar solicitação ao suporte}{%
\begin{elementos}
\item \textbf{Assunto.} Em Ajuda e suporte, resuma a necessidade.
\item \textbf{Como podemos ajudar?} Descreva a dúvida, sugestão ou problema. O campo aceita até 5.000 caracteres.
\item \textbf{Enviar solicitação.} Envie a mensagem à equipe de suporte.
\end{elementos}
}

\passo[0.89560/0.22373/0.85714/0.16949,0.38462/0.51977/0.38462/0.51977]{completo/66-solicitacoes-suporte}{Consultar solicitações de suporte}{
Com perfil Administrador, abra Suporte no menu lateral para acessar Solicitações de suporte.\par
\begin{elementos}
\item \textbf{Buscar solicitações.} Use nome, e-mail ou assunto para localizar um chamado.
\item \textbf{Caixa de entrada.} Consulte as mensagens recebidas. Se a lista estiver vazia, limpe a busca para conferir todas as solicitações.
\end{elementos}
}

\Needspace{160pt}

\subsection*{Campos e regras}

\begingroup
\fontsize{9}{12}\selectfont
\setlength{\LTpre}{4pt}
\setlength{\LTpost}{8pt}
\begin{longtable}{@{}L{3.389cm}L{1.916cm}L{2.652cm}L{6.778cm}@{}}
\toprule
\guiacabtabela \thead{Campo} & \thead{Obrigatório} & \thead{Formato} & \thead{Descrição e regra de preenchimento} \\
\midrule
\endfirsthead
\toprule
\guiacabtabela \thead{Campo} & \thead{Obrigatório} & \thead{Formato} & \thead{Descrição e regra de preenchimento} \\
\midrule
\endhead
\bottomrule
\endlastfoot
\rowcolor{white} \textbf{Assunto} & Sim & Texto curto & Resumo do tema da solicitação. \\
\rowcolor{guiaareia!70} \textbf{Como podemos ajudar?} & Sim & Texto longo (até 5.000 caracteres) & Explicação detalhada da dúvida, sugestão ou ocorrência. \\
\rowcolor{white} \textbf{Buscar solicitações} & Não & Texto livre & Filtra a caixa de entrada do administrador por nome, e-mail ou assunto. \\
\end{longtable}
\endgroup

\subsection*{Mensagens e solução de problemas}

\begin{remark}
\textbf{Nenhuma solicitação encontrada — As mensagens enviadas aparecerão nesta lista.}\par
\textbf{Causa provável:} Ainda não há chamados abertos ou o texto digitado na busca não encontrou correspondência.\par
\textbf{Como resolver:} Limpe o campo \textbf{Buscar solicitações} para ver toda a caixa de entrada.
\end{remark}

\section{Cadastro e Gestão de Usuários}

\subsection*{Objetivo}

Cadastrar novos colaboradores no sistema, definir seus cargos de acesso (\tech{Analista}, \tech{Gestor} ou \tech{Admin}), ativar ou desativar contas e manter os dados da equipe atualizados. A gestão de contas exige permissão administrativa de \textbf{superusuário}.\par

\subsection*{Antes de começar}

\begin{itemize}
\item Estar conectado com permissão de \textbf{superusuário}.
\end{itemize}

\passo[0.45109/0.54545/0.45109/0.54545,0.92174/0.10227/0.90761/0.06818,0.74130/0.50341/0.74457/0.44318,0.88587/0.50341/0.85326/0.56818,0.92065/0.50341/0.96196/0.56818]{completo/67-usuarios}{Consultar os usuários}{%
\begin{elementos}
\item \textbf{Tabela.} Em Usuários, confira nome, usuário, e-mail, gestor direto e cargo de cada conta.
\item \textbf{Novo usuário.} Abra o formulário de cadastro, detalhado no passo~\ref{passo:57}.
\item \textbf{Ativo.} Confira o status da conta na chave correspondente.
\item \textbf{Editar.} Use o lápis da linha para alterar os dados.
\item \textbf{Excluir.} Use a lixeira da conta quando a remoção for necessária e autorizada.
\end{elementos}
}

\passo[0.50000/0.18171/0.13889/0.13415,0.50000/0.32561/0.13889/0.28049,0.50000/0.46707/0.13889/0.42683,0.50000/0.61098/0.13889/0.57317,0.50000/0.75488/0.13889/0.70732,0.82778/0.91098/0.75000/0.85366]{completo/68-novo-usuario}{Cadastrar novo usuário}{%
\begin{elementos}
\item \textbf{Nome.} Informe o nome de exibição do colaborador, de 2 a 30 caracteres.
\item \textbf{Usuário.} Informe de 2 a 20 caracteres, apenas letras minúsculas e números, sem espaços ou símbolos.
\item \textbf{E-mail.} Informe o e-mail corporativo, com até 50 caracteres.
\item \textbf{Cargo.} Escolha Analista, Gestor ou Admin.
\item \textbf{Senha inicial.} Use ao menos 12 caracteres. A senha não deve conter nome, usuário ou e-mail do colaborador.
\item \textbf{Criar usuário.} Conclua o cadastro.
\end{elementos}
}

\Needspace{160pt}

\subsection*{Campos e regras}

\begingroup
\fontsize{9}{12}\selectfont
\setlength{\LTpre}{4pt}
\setlength{\LTpost}{8pt}
\begin{longtable}{@{}L{3.389cm}L{1.916cm}L{2.652cm}L{6.778cm}@{}}
\toprule
\guiacabtabela \thead{Campo} & \thead{Obrigatório} & \thead{Formato} & \thead{Descrição e regra de preenchimento} \\
\midrule
\endfirsthead
\toprule
\guiacabtabela \thead{Campo} & \thead{Obrigatório} & \thead{Formato} & \thead{Descrição e regra de preenchimento} \\
\midrule
\endhead
\bottomrule
\endlastfoot
\rowcolor{white} \textbf{Nome} & Sim & Texto (2 a 30 caracteres) & Nome do colaborador exibido nas telas, tarefas e relatórios. \\
\rowcolor{guiaareia!70} \textbf{Usuário} & Sim & Letras minúsculas e números (2 a 20 caracteres) & Identificação usada no login; não aceita espaços, acentos ou letras maiúsculas. \\
\rowcolor{white} \textbf{E-mail} & Sim & Endereço de e-mail (até 50 caracteres) & E-mail corporativo vinculado ao colaborador. \\
\rowcolor{guiaareia!70} \textbf{Cargo} & Sim & Campo de seleção (\tech{Analista}, \tech{Gestor}, \tech{Admin}) & Define quais telas e ações o colaborador poderá acessar. \\
\rowcolor{white} \textbf{Senha inicial} & Sim & Texto oculto (12 a 72 caracteres) & Deve ter no mínimo 12 caracteres, ser difícil de adivinhar e não conter o nome, usuário ou e-mail da pessoa. \\
\rowcolor{guiaareia!70} \textbf{STATUS (``Ativo'')} & Sim & Chave liga/desliga & Permite desativar ou reativar imediatamente o acesso do colaborador. \\
\end{longtable}
\endgroup

\subsection*{Mensagens e solução de problemas}

\begin{remark}
\textbf{Aviso de senha muito simples ou fácil de adivinhar}\par
\textbf{Causa provável:} A senha informada tem menos de 12 caracteres, usa uma sequência muito comum ou contém partes do nome, usuário ou e-mail do colaborador.\par
\textbf{Como resolver:} Crie uma senha com pelo menos 12 caracteres sem repetir o nome ou login do usuário.
\end{remark}

\begin{remark}
\textbf{Erro ao validar o campo Usuário ou Nome}\par
\textbf{Causa provável:} Foram digitadas letras maiúsculas ou espaços no campo \textbf{Usuário}, ou o \textbf{Nome} passou de 30 caracteres.\par
\textbf{Como resolver:} Use apenas letras minúsculas e números no campo \textbf{Usuário} e abrevie nomes intermediários se ultrapassar 30 caracteres.
\end{remark}

\Needspace{190pt}
\section{Uso da Plataforma e Configuração de Tempos Manuais}

\subsection*{Objetivo}

Acompanhar o tempo economizado pela equipe com o uso do sistema Hanaro (no dia, no mês e no ano), consultar a quantidade de acessos diários e ajustar os tempos médios que cada atividade levava antes da automação. Exclusivo para \textbf{Administradores}.\par

\subsection*{Antes de começar}

\begin{itemize}
\item Estar conectado com perfil de \textbf{Administrador}.
\end{itemize}

\passo[0.50543/0.24309/0.50543/0.24309,0.49457/0.45635/0.49457/0.45635,0.50543/0.74033/0.50543/0.74033,0.92935/0.54144/0.89674/0.58564]{completo/69-uso-plataforma}{Consultar o uso da plataforma}{%
\begin{elementos}
\item \textbf{Tempo economizado.} Em Uso da plataforma, compare Hoje, Este mês e Este ano, observando tempo atual e manual estimado.
\item \textbf{Utilização.} Confira Acessos, Usuários ativos, Coletas concluídas e Relatórios gerados.
\item \textbf{Acessos por dia.} Acompanhe a evolução de utilização.
\item \textbf{Período.} Alterne a exibição entre os últimos 30 e 90 dias.
\end{elementos}
}

\passo[0.50000/0.35397/0.13889/0.30159,0.50000/0.54444/0.13889/0.49206,0.50000/0.73651/0.13889/0.68254,0.82222/0.88413/0.75000/0.95238]{completo/70-tempos-manuais}{Configurar os tempos manuais}{
Em Uso da plataforma, clique em Configurar para abrir Configurar tempos manuais.\par%
\begin{elementos}
\item \textbf{Preparação da coleta.} Em Configurar tempos manuais, informe os minutos por tarefa, como 90 minutos.
\item \textbf{Apresentação dos dados.} Informe a estimativa manual, como 45 minutos.
\item \textbf{Criação de relatório.} Informe a estimativa para elaborar um relatório, como 120 minutos.
\item \textbf{Salvar tempos.} Grave as estimativas usadas no cálculo de tempo economizado.
\end{elementos}
}

\Needspace{160pt}

\subsection*{Campos e regras}

\begingroup
\fontsize{9}{12}\selectfont
\setlength{\LTpre}{4pt}
\setlength{\LTpost}{8pt}
\begin{longtable}{@{}L{3.389cm}L{1.916cm}L{2.652cm}L{6.778cm}@{}}
\toprule
\guiacabtabela \thead{Campo} & \thead{Obrigatório} & \thead{Formato} & \thead{Descrição e regra de preenchimento} \\
\midrule
\endfirsthead
\toprule
\guiacabtabela \thead{Campo} & \thead{Obrigatório} & \thead{Formato} & \thead{Descrição e regra de preenchimento} \\
\midrule
\endhead
\bottomrule
\endlastfoot
\rowcolor{white} \textbf{Preparação da coleta} & Sim & Número inteiro (\tech{minutos por tarefa}) & Tempo que era gasto manualmente para extrair e organizar cada carga de dados antes da automação. \\
\rowcolor{guiaareia!70} \textbf{Apresentação dos dados} & Sim & Número inteiro (\tech{minutos por tarefa}) & Tempo que era gasto manualmente para atualizar planilhas e gráficos de indicadores. \\
\rowcolor{white} \textbf{Criação de relatório} & Sim & Número inteiro (\tech{minutos por tarefa}) & Tempo que era gasto manualmente para montar e formatar cada relatório executivo. \\
\rowcolor{guiaareia!70} \textbf{Período do gráfico (``30 dias'' / ``90 dias'')} & Não & Botões de opção & Muda o período exibido no gráfico de \textbf{Acessos por dia}. \\
\end{longtable}
\endgroup

\subsection*{Mensagens e solução de problemas}

\begin{remark}
\textbf{Tempo economizado marcando ``0min''}\par
\textbf{Causa provável:} Não houve coletas ou relatórios gerados no período consultado (como \tech{Hoje}) ou os valores em \textbf{Configurar tempos manuais} estão zerados.\par
\textbf{Como resolver:} Clique em \textbf{\faCog{} Configurar} no canto superior direito e verifique se os minutos por tarefa estão preenchidos.
\end{remark}

\fimdecapitulo

\chapter{Assistente de IA e Consulta}\label{cap:assistente-uso}

\section{Disponibilidade e uso das propostas}

O Assistente prepara propostas a partir dos dados do Hanaro e de análises já
aprovadas. Cabe ao Analista ou Gestor conferir os fatos e decidir o que aproveitar.
Uma sugestão de causa deve ser tratada como hipótese até a verificação na linha.

\begin{remark}[info]
Assistente, Consulta, Memória, Desempenho e tradução dependem dos recursos
disponibilizados no ambiente e das permissões da conta. Se uma opção não aparecer,
confirme sua disponibilidade com o responsável pelo sistema. As figuras deste
capítulo são \textbf{esquemas ilustrativos de referência}; os rótulos e a disposição
dos controles podem variar. Elas não representam capturas reais de uma sessão.
\end{remark}
\begingroup
\fontsize{9}{12}\selectfont
\begin{longtable}{@{}L{3.1cm}L{12.1cm}@{}}
\toprule
\guiacabtabela \thead{Função} & \thead{Quem utiliza, quando disponível} \\
\midrule
\endfirsthead
\toprule
\guiacabtabela \thead{Função} & \thead{Quem utiliza, quando disponível} \\
\midrule
\endhead
\bottomrule
\endlastfoot
\rowcolor{white} Pré-análise e Consulta & Analista e Gestor, nas opções liberadas para o ambiente. \\
\rowcolor{guiaareia!70} Memória & Gestor e administrador autorizado pela manutenção da memória. \\
\rowcolor{white} Desempenho & Perfis de gestão autorizados a consultar os indicadores. \\
\rowcolor{guiaareia!70} Tradução & Usuários com acesso à emissão do relatório e ao recurso de tradução. \\
\end{longtable}
\endgroup

Uma pré-análise pode ser preparada a partir de uma regra de alerta configurada
para solicitar apoio do Assistente. A proposta pode seguir a rota de
\textbf{Recorrência}, \textbf{Caso parecido} ou \textbf{Sem precedente}. Confira a
rota antes de usar as sugestões: um caso semelhante não comprova a mesma causa.

\Needspace{330pt}
\section{Abrir e conferir a pré-análise}

\passo[0.07000/0.14911/0.08000/0.22367,0.35000/0.09586/0.15500/0.09586,0.35000/0.20237/0.15500/0.25562,0.35000/0.38343/0.15500/0.44734]{ia/01-assistente-lista}{Abrir a lista do Assistente}{%
\begin{elementos}
\item \textbf{Assistente.} Abra o menu para consultar as pré-análises. O contador indica as propostas prontas que aguardam decisão.
\item \textbf{Seções.} Use Pré-análises, Consulta, Memória ou Desempenho conforme os acessos da conta.
\item \textbf{Abas e filtros.} Escolha Aguardando decisão, Decididas ou Todas e refine a lista por gatilho, rota, linha ou período.
\item \textbf{Pré-análise.} Confira item, linha, valor, gatilho e rota. Abra a proposta desejada para ler os detalhes.
\end{elementos}
}

\passo[0.35000/0.12781/0.15500/0.17041,0.35000/0.42604/0.15500/0.47929,0.35000/0.70296/0.15500/0.75621,0.70000/0.70296/0.61000/0.77751]{ia/02-preanalise-leitura}{Conferir fatos, sugestões e fontes}{%
\begin{elementos}
\item \textbf{Rota e fatos.} Leia O que aconteceu e confira item, linha, ocorrências, período e valores com os registros do caso.
\item \textbf{Proposta de análise.} Revise título, descrição, tipo de scrap, categoria 4M, hipótese de causa e contenção.
\item \textbf{Casos parecidos.} Abra as análises citadas e confira se a situação é aplicável. Use Não é parecido quando a comparação não fizer sentido.
\item \textbf{Justificativa e fontes.} Leia por que cada campo foi sugerido e confira a origem dos dados antes de aceitar a proposta.
\end{elementos}
}

\Needspace{330pt}
\section{Decidir sobre a proposta}

\passo[0.35000/0.12781/0.15500/0.17041,0.35000/0.38343/0.15500/0.42604,0.35000/0.63905/0.15500/0.68166,0.35000/0.85207/0.15500/0.90533]{ia/03-preanalise-decisao}{Aceitar, editar ou rejeitar sugestões}{%
\begin{elementos}
\item \textbf{Campos propostos.} Decida entre Aceitar, Editar e Rejeitar para cada campo. Use Desfazer para rever a decisão e confira o conteúdo antes de Aceitar tudo.
\item \textbf{Tarefas propostas.} Revise separadamente as ações sugeridas. Elas só serão criadas quando forem aceitas e vinculadas a um relatório.
\item \textbf{Relatório e finalização.} Se quiser finalizar a análise e criar o relatório, marque a opção correspondente e informe os dados solicitados. Sem essa opção, a análise permanece em rascunho.
\item \textbf{Aplicar decisões.} Grave as escolhas. Dispensar encerra a proposta sem aproveitá-la e solicita um motivo. Quando aparecer a avaliação A pré-análise ajudou?, registre sua percepção.
\end{elementos}
}

\begin{remark}
Se o sistema informar que existe uma versão mais nova da pré-análise ou que ela
mudou durante a revisão, use \textbf{Carregar a versão nova} e confira novamente
os dados antes de aplicar. Em propostas para um grupo de ocorrências, confira
quais análises e qual relatório serão criados; as análises já existentes são mantidas.
\end{remark}
\begingroup
\fontsize{9}{12}\selectfont
\begin{longtable}{@{}L{3.1cm}L{12.1cm}@{}}
\toprule
\guiacabtabela \thead{Situação} & \thead{Significado} \\
\midrule
\endfirsthead
\toprule
\guiacabtabela \thead{Situação} & \thead{Significado} \\
\midrule
\endhead
\bottomrule
\endlastfoot
\rowcolor{white} Aguardando decisão & A proposta está pronta e ainda depende da conferência de uma pessoa. \\
\rowcolor{guiaareia!70} Aplicada & As decisões foram gravadas e os itens aceitos foram criados. \\
\rowcolor{white} Aplicada em parte & Parte das decisões foi gravada; confira os avisos para identificar o que ainda precisa de tratamento. \\
\rowcolor{guiaareia!70} Dispensada & A proposta foi descartada, com o motivo informado. \\
\rowcolor{white} Substituída & Uma proposta mais nova passou a ser a referência; a anterior fica no histórico. \\
\end{longtable}
\endgroup

\Needspace{330pt}
\section{Consultar casos já tratados}

\passo[0.78000/0.09586/0.65000/0.15976,0.35000/0.83077/0.15500/0.88402,0.35000/0.23432/0.15500/0.28757,0.35000/0.52189/0.15500/0.58580]{ia/04-consulta}{Perguntar na Consulta}{%
\begin{elementos}
\item \textbf{Conversas.} Use Nova conversa para iniciar uma consulta ou Histórico para retomar uma conversa. A opção de apagar remove a conversa escolhida.
\item \textbf{Pergunta.} Descreva o que procura em texto livre, por exemplo como foi tratado um defeito em determinada linha, e clique em Perguntar.
\item \textbf{Resposta.} Leia a orientação e confira se ela responde à situação descrita. A Consulta utiliza análises aprovadas disponíveis na memória.
\item \textbf{Fontes.} Use Abrir análise para conferir os registros citados. Registre Útil ou Não útil para avaliar a resposta.
\end{elementos}
}

Se não houver resultados, tente outras palavras, o código do item ou a linha.
A ausência de casos semelhantes não confirma uma causa nem dispensa a investigação.

\clearpage
\section{Memória e desempenho}

\passo[0.35000/0.12781/0.15500/0.18107,0.35000/0.30888/0.15500/0.36213,0.35000/0.42604/0.15500/0.50059,0.35000/0.85207/0.15500/0.90533]{ia/05-memoria}{Acompanhar a memória do Assistente}{%
\begin{elementos}
\item \textbf{Indicadores.} Confira as análises aprovadas, os itens na memória, os que aguardam envio, os retirados e eventuais falhas.
\item \textbf{Busca e situação.} Localize uma análise pelo título ou item e filtre sua situação na memória.
\item \textbf{Análises aprovadas.} Confira item, linha e data. Retirar da memória suspende o uso do caso como exemplo; Colocar de volta restaura sua participação.
\item \textbf{Reenviar tudo.} Use quando houver orientação para reconstruir ou atualizar a memória. O envio acontece em lotes; acompanhe a fila e as falhas após a solicitação.
\end{elementos}
}

\Needspace{100pt}
Na aba \textbf{Desempenho}, selecione o período e acompanhe o total de propostas
decididas, a parcela aceita sem alteração, o tempo de preparação, os campos
rejeitados e seus motivos e as avaliações registradas. Esses indicadores ajudam
a revisar o uso do Assistente; o aceite de uma proposta não comprova o resultado
da ação na fábrica.

\section{Tradução e emissão do relatório}

A emissão e o acompanhamento da tradução estão no passo~\ref{passo:45}
(\pagina{passo:45}); o formato de apresentação está no passo~\ref{passo:46}.
Use essas orientações para conferir a fila, o andamento, os trechos mantidos em
português e as alternativas em caso de falha, sem repetir o procedimento de exportação.

\fimdecapitulo


% ---------- FECHAMENTO ----------
\guiapagina{Fontes, créditos e referência de versão}{%
    Esta edição integra o guia rápido original e as orientações operacionais da
    Documentação Técnica do Hanaro, produzidos pela equipe do Polo INOVA.
    O histórico do repositório registra contribuições de Messyas e de
    André Filipe Aloise (\texttt{aaloise}); as capturas reais acrescentadas à
    Apostila e às Configurações foram incorporadas ao repositório pelo PR~1.

    As orientações de Configurações e do Assistente foram adaptadas para este
    manual. Os esquemas do Assistente são os do documento existente e estão
    identificados como ilustrações. As capturas originais foram preservadas.

    A referência funcional desta edição é a versão apresentada nas capturas e
    nos manuais de outubro e novembro de 2026. Os recursos de IA e alguns
    controles dependem da versão, do ambiente e das permissões da conta.
    Os capítulos 1 e 2 mantêm o guia rápido da versão-base; suas indicações
    de escopo e tempo de leitura se referem ao uso básico.

    A arquitetura, os fluxos internos de IA e o modelo de dados permanecem na
    \emph{Documentação Técnica do Hanaro}, em
    \href{https://github.com/Messyas/templates-ax-academy-latex/tree/main/modelos/Guia-Usuario-Hanaro}{modelos/Guia-Usuario-Hanaro}.

    Fontes do conteúdo: guia rápido, versão-base \texttt{20bda03}; Documentação
    Técnica e capturas do repositório até \texttt{de705b5}; consolidação do guia
    operacional na branch \texttt{doc/guia-hanaro}.
}

\guiacontracapa

\end{document}
